% !TEX root = main.tex
% ============================================================
%  DKLS23 学习笔记 · 单栏论文模板
%  article + xeCJK: 英文论文版式, 中文宋体正文 / 黑体标题
%  全部宏包 TeX Live 自带; 无页眉 (页脚中央页码)
%
%  每篇笔记一份 PDF. 标题由片段的 "% title: ..." 注释自动注入 (template/wrap.py),
%  无作者 / 日期 / 摘要 / 关键词. 唯一手动开关是 \input 行
%  (文件里唯一一个 \input, make 生成包装文件时自动替换).
%
%  编译: make                    # 全量 14 篇 (任何片段改动都会自动全量重编)
%        make pdf/06-rvole.pdf   # 只编译一篇
%  (产物在 pdf/)
% ============================================================
% article 类只支持 10/11/12pt, 14pt 要换 extarticle (article 的同款克隆, TeX Live 自带)
\documentclass[14pt,fleqn]{extarticle}
% 四边均匀砍掉 1/3: 2.6cm → 1.73cm
\usepackage[top=1.73cm,left=1.73cm,right=1.73cm,bottom=1.73cm]{geometry}
% 数学正文使用 New Computer Modern Math；黑板粗体使用经典的有衬线 AMS 字形。
% unicode-math 必须最后加载（amsmath/amssymb 之后）。
\usepackage{amsmath,amssymb,mathtools}
% 单栏协议：自动横线、缩进辅助线与显式作用域。
\makeatletter
\providecommand\input@path{}
\g@addto@macro\input@path{{../template/}{template/}}
\makeatother
\usepackage{ruledprotocol}
\usepackage{unicode-math}
\setmathfont{NewCMMath-Regular.otf}
% 粗体数学版本（\boldsymbol / \mathversion{bold} 依赖）使用配套的 NewCM Math Bold。
\setmathfont[version=bold]{NewCMMath-Bold.otf}
\usepackage[bb=ams]{mathalpha}
\newcommand{\bit}{\ensuremath{\left\{0,1\right\}}}
\newcommand{\bits}[1]{\ensuremath{\bit^{#1}}}
\newcommand{\true}{\ensuremath{\mathrm{True}}}
\newcommand{\false}{\ensuremath{\mathrm{False}}}
\newcommand{\hash}[1]{\ensuremath{\mathrm{Hash}\left(#1\right)}}
\newcommand{\prf}[1]{\ensuremath{\mathrm{PRF}\left(#1\right)}}
\newcommand{\uniform}[1]{\ensuremath{\mathrm{Unif}\left(#1\right)}}
\newcommand{\enc}[1]{\ensuremath{\mathrm{Enc}\left(#1\right)}}
\newcommand{\dec}[1]{\ensuremath{\mathrm{Dec}\left(#1\right)}}
\newcommand{\call}[2]{\ensuremath{\mathbf{#1}\left(#2\right)}}
% 伪代码 (《算法导论》体例): Cormen 自带的 clrscode3e, \begin{codebox} + \li 编号行。
\usepackage{clrscode3e}
\usepackage{longtable,booktabs,array}
\usepackage{enumitem}
\setlist{nosep}

% 中文：正文思源宋体，标题思源黑体。
\usepackage{xeCJK}
% 正文西文同样使用思源宋体；思源无斜体面，以楷体替代。
\setmainfont{NotoSerifSC-Regular.ttf}[
  BoldFont=NotoSerifSC-Bold.ttf,
  ItalicFont=FandolKai-Regular.otf,
]
\setCJKmainfont{NotoSerifSC-Regular.ttf}[
  BoldFont=NotoSerifSC-Bold.ttf,
  ItalicFont=FandolKai-Regular.otf,
]
\setCJKsansfont{NotoSansSC-Regular.ttf}[
  BoldFont=NotoSansSC-Bold.ttf,
]
\setCJKmonofont{FandolFang-Regular.otf}
\newCJKfontfamily\heiti{NotoSansSC-Medium.ttf}[BoldFont=NotoSansSC-Bold.ttf]
\newfontfamily\heitiLatin{NotoSansSC-Medium.ttf}[BoldFont=NotoSansSC-Bold.ttf]

% 行间公式左对齐，按正文中文字体实测两个汉字的宽度作为缩进。
\AtBeginDocument{\settowidth{\mathindent}{汉字}}

% --- 章节标题黑体 (中英一致无衬线), 论文观感 ---
% 标题字体本身已是 SemiBold, 不加 \bfseries (无衬线再粗就笨重)
% 编号显示加英文句点结尾 ("1." "1.1."), 只改显示不改 \thesection → \ref 引用仍是干净的 "1.1"
\usepackage{titlesec}
\titleformat{\section}{\Large\heitiLatin\heiti}{\thesection.}{1em}{}
\titleformat{\subsection}{\large\heitiLatin\heiti}{\thesubsection.}{1em}{}
\titleformat{\subsubsection}{\normalsize\heitiLatin\heiti}{\thesubsubsection.}{1em}{}

% --- 笔记排版: 段间空行, 段首不缩进 ---
% 行距 1.3 倍 (ctex/ustcthesis 标准: 1.2 基线 × 1.3 = 1.56 倍字号; article 默认只有 1.2)
\linespread{1.3}
\setlength{\parindent}{0pt}
% 段距 1em: 0.6em 在 1.3 行距下与行距混成一片, 没有段落区分度
\setlength{\parskip}{1em}

% 补丁使用可跨页的圆角虚线框，首段自动以 🛡️ 开头。
\usepackage[most]{tcolorbox}
\usepackage{twemojis}
\newenvironment{patchnote}[1]{%
  \begin{tcolorbox}[
    enhanced jigsaw, breakable,
    colback=white, frame hidden,
    borderline={0.6pt}{0pt}{black!65,dashed},
    arc=3mm, boxrule=0.6pt,
    % 分页处用水平虚线闭合，直接连接左右边框，不加圆角。
    overlay first={
      \draw[black!65,dashed,line width=0.6pt]
        ([xshift=0.3pt]frame.south west) -- ([xshift=-0.3pt]frame.south east);
    },
    overlay middle={
      \draw[black!65,dashed,line width=0.6pt]
        ([xshift=0.3pt]frame.north west) -- ([xshift=-0.3pt]frame.north east);
      \draw[black!65,dashed,line width=0.6pt]
        ([xshift=0.3pt]frame.south west) -- ([xshift=-0.3pt]frame.south east);
    },
    overlay last={
      \draw[black!65,dashed,line width=0.6pt]
        ([xshift=0.3pt]frame.north west) -- ([xshift=-0.3pt]frame.north east);
    },
    boxsep=1mm, left=2mm, right=2mm, top=2mm, bottom=2mm,
    grow to left by=3.6mm, grow to right by=3.6mm,
    before skip=1em, after skip=1em,
    before upper={\setlength{\parindent}{0pt}\setlength{\parskip}{1em}}
  ]
  {\heitiLatin\heiti\texttwemoji{shield}\enspace #1}\par\nopagebreak
}{%
  \end{tcolorbox}
}

% --- 页眉: 无. article 默认 pagestyle=plain, 只有页脚中央页码 ---
\pagestyle{plain}

% --- 跨笔记引用: 其它 14 篇的 .aux 提供标签, 链接跳到 pdf/ 里的 PDF ---
%     make 全量编译时自动刷新; 单篇编译注意 stale aux (见 Makefile 注释)
\usepackage{xr-hyper}
\externaldocument{00-mta-baseot}
\externaldocument{01-iknp03}
\externaldocument{02-kos15}
\externaldocument{03-endemic-ot}
\externaldocument{04-pprf}
\externaldocument{05-softspoken}
\externaldocument{06-rvole}
\externaldocument{07-orchestration}
\externaldocument{dkls23}
\externaldocument{misc-f2k}
\externaldocument{misc-fiat-shamir}
\externaldocument{misc-gadget}
\externaldocument{misc-rerand}
\externaldocument{uc-framework}

\usepackage{hyperref}
\hypersetup{colorlinks=true,linkcolor=blue,urlcolor=blue,citecolor=blue}
\usepackage{pdfcomment}

% --- 智能引用 cleveref: \cref{a,b,c} 自动连多个引用 + 类型前缀 ---
%     必须最后加载 (hyperref 之后), \cref 才能生成可点击链接。
\usepackage{cleveref}
\crefname{equation}{公式}{公式}
\Crefname{equation}{公式}{公式}
\crefname{section}{节}{节}
\Crefname{section}{节}{节}
\crefname{subsection}{节}{节}
\Crefname{subsection}{节}{节}
% 多引用连接词 (默认英文 and/,): 两个→" 和 ", 三个→"、... 和 ..."
% 注意: 这些命令 cleveref 在 \AtBeginDocument 才定义, 这里要用 \newcommand 预定义
% (cleveref 检测到已定义就会保留, \renewcommand 会因"尚未定义"报错)
\newcommand{\crefpairconjunction}{ 和 }
\newcommand{\crefmiddleconjunction}{、}
\newcommand{\creflastconjunction}{ 和 }
\newcommand{\crefrangeconjunction}{ 至 }

% ==================== 每篇笔记改这里 ====================
% 标题由 Makefile 从片段的 "% title: ..." 注释注入到 \doctitle (见 template/wrap.py)。
% 中英文都使用思源黑体，\mdseries 抵消默认粗体。
% 只显示标题, 无作者 / 日期 / 摘要 / 关键词。
\providecommand{\doctitle}{未命名笔记}

\begin{document}
\begin{center}
  {\LARGE\heitiLatin\heiti\mdseries\doctitle\par}
\end{center}
\vspace{1em}

% >>> 要编译的笔记片段: 改这一行 <<<
% !TEX root = 00-mta-baseot.tex
% title: DKLS23 实施方式

\section{前置工作}

\subsection{朴素 OT-MtA}

\begin{protocol}[
  name=EcdsaSign,
  params={m,x \mid n,G},
  returns={r,s}
]
  k \coloneq \uniform{\mathbb{Z}_n^*} \\
  R \coloneq kG \\
  r \coloneq R.x \bmod n \pccomment{点 $R$ 的横坐标} \\
  \label{ecdsa-s}\tag{ecdsa-s}
  \boldsymbol{s} \coloneq k^{-1}(m+rx) \bmod n \\
\end{protocol}

为 \eqref{ecdsa-s} 引入乘项 $\gamma \ne 0 \pmod n$，则有
\begin{equation*}
  s = (k\gamma)^{-1}(m\gamma + rx\gamma) \bmod n.
\end{equation*}

朴素 MtA 协议：

\begin{protocol}[
  name=NaiveMtA,
  params={\party{S}: x_a,~ \party{R}: x_b \mid n},
  returns={\party{S}: y_a,~ \party{R}: y_b},
  st={y_a+y_b \equiv x_a x_b \pmod n}
]
  \pcparty{S}
  \pcfor t \in [0, \ell); \pcdo \\
  r_t \coloneq \uniform{\mathbb{Z}_n^*} \\
  m_{t,0} \coloneq r_t,~ m_{t,1} \coloneq r_t + x_a 2^t \bmod n \\
  \pcendfor{t} \\

  \pcparty{R}\sum_{t=0}^{\ell-1}d_t 2^t \coloneq x_b \\

  \pcparty{S,R}\pcfor t \in [0, \ell); \pcdo \\
    \party{S}: \varnothing,~ \party{R}: m_{t,d_t} \coloneq
    \mathrm{NaiveOT}\left(\party{S}: (m_{t,0}, m_{t,1}),~ \party{R}: d_t\right) \\
  \pcendfor{t} \\

  \pcparty{S}y_a \coloneq -\sum_t m_{t,0} \bmod n \\
  \pcparty{R}y_b \coloneq \sum_t m_{t, d_t} \bmod n \\
\end{protocol}

朴素二选一 OT 协议：

\begin{protocol}[
  name=NaiveOT,
  params={\party{S}: (m_0, m_1),~ \party{R}: d\in\bit},
  returns={\party{S}: \emptyset,~ \party{R}: m_d}
]
  \pcparty{S}
  \alpha \coloneq \uniform{\mathbb{Z}_n^*},~ A \coloneq \alpha G \\
  \pcsend{A}{\party{R}} \\

  \pcparty{R}
  \beta \coloneq \uniform{\mathbb{Z}_n^*},~ B \coloneq \beta G + dA \\
  \pcsend{B}{\party{S}} \\

  \pcparty{S}
  K_0 \coloneq \hash{\alpha B},~ K_1 \coloneq \hash{\alpha (B-A)} \\
  C_0 \coloneq \enc{K_0, m_0},~ C_1 \coloneq \enc{K_1, m_1} \\
  \pcsend{C_0, C_1}{\party{R}} \\
  \pcparty{R}
  K_d \coloneq \hash{\beta A} \\
  m_d \coloneq \dec{K_d, C_d} \\
  \pcreturn
\end{protocol}
\bigskip

\subsection{IKNP03}

制备持久化种子，用于后续 OT 协议。

\begin{protocol}[
  name=BaseOTSetup,
  params={\kappa},
  returns={\party{S}: (s, k^s),~ \party{R}: (k^0, k^1)}
]
  \pcparty{S,R}\pcfor i \in [0,\kappa); \pcdo \\
  \pcparty{R}
  k_i^0 \coloneq \uniform{\bits{\kappa}},~
  k_i^1 \coloneq \uniform{\bits{\kappa}} \\
  \pcparty{S}s_i \coloneq \uniform{\bit} \\
  \pcparty{S,R}
  \party{R}: \varnothing,~ \party{S}: k_i^{s_i} \coloneq
  \call{NaiveOT}{\party{R}: (k_i^0,k_i^1),~ \party{S}: s_i}
  \pccomment{注意角色互换}
  \\
  \pcendfor{i} \\
  \pcparty{S}
  s \coloneq \bigl(s_i \mid i\in[0,\kappa)\bigr),~
  k^s \coloneq \bigl(k_i^{s_i} \mid i\in[0,\kappa)\bigr) \\
  \pcparty{R}
  k^0 \coloneq \bigl(k_i^0 \mid i\in[0,\kappa)\bigr),~
  k^1 \coloneq \bigl(k_i^1 \mid i\in[0,\kappa)\bigr) \\
  \pcreturn
\end{protocol}
\bigskip

子协议 IknpExpand 意在用 $\kappa$ 个（对）种子，生成 $m$ 个（对）密钥原像。
Sender 获得成对的密钥原像 $q_{*,j}$ 与 $q_{*,j}\oplus s$；Receiver 获得所选密钥原像 $t^0$。
\pdfcomment[
  author={wojtek},
  subject={Test Comment}
]{
  Whatever.
}

\begin{protocol}[
  name=IknpExpand,
  params={\party{S}: (s,k^s),~ \party{R}: (k^0,k^1,b) \mid \mathtt{sid}, \mathtt{tag}, L},
  returns={\party{S}: (q,u),~ \party{R}: (t^0,u)},
  st={\forall j\in[0,L),~ q_{*,j}=t^0_{*,j}\oplus(b_j\odot s)},
  % layout=multiline
]
  \pcparty{R}\pcfor i \in [0,\kappa); \pcdo \\
  t^0_{i,*} \coloneq \prf{\mathtt{sid}, \mathtt{tag},k_i^0} \in \bits{L} \\
  t^1_{i,*} \coloneq \prf{\mathtt{sid}, \mathtt{tag},k_i^1} \in \bits{L} \\
  u_{i,*} \coloneq t^0_{i,*} \oplus t^1_{i,*} \oplus b \\
  \pcendfor{i} \\
  \pcsend{u}{\party{S}} \\
  \pcparty{S}\pcfor i \in [0,\kappa); \pcdo \\
  t_{i,*} \coloneq \prf{\mathtt{sid}, \mathtt{tag},k_i^{s_i}} \in \bits{L} \\
  q_{i,*} \coloneq t_{i,*} \oplus (s_i \odot u_{i,*}) \\
  \pcendfor{i} \\
  \pcreturn
  \pcnote{
    操作符 $\oplus$ 是按位异或，$\odot$ 是按位与；
    矩阵 $q,u,t^0$ 的形状均为 $\kappa\times L$。}
\end{protocol}

双方算出密钥，传输二选一消息。

\begin{protocol}[
  name=PayloadOT,
  params={\party{S}: (q,s,M^0,M^1),~ \party{R}: (t^0,b) \mid \mathtt{sid}, \mathtt{tag},m},
  returns={\party{S}: \varnothing,~ \party{R}: M^b},
  layout=multiline
]
  \pcparty{S}\pcfor j \in [0,m); \pcdo \\
  K_j^0 \coloneq \hash{\mathtt{sid}, \mathtt{tag},j,q_{*,j}} \\
  K_j^1 \coloneq \hash{\mathtt{sid}, \mathtt{tag},j,q_{*,j}\oplus s} \\
  C_j^0 \coloneq \enc{K_j^0,M_j^0},~
  C_j^1 \coloneq \enc{K_j^1,M_j^1} \\
  \pcsend{j,C_j^0,C_j^1}{\party{R}} \\
  \pcendfor{j} \\
  \pcparty{R}\pcfor j \in [0,m); \pcdo \\
  K_j^{b_j} \coloneq \hash{\mathtt{sid}, \mathtt{tag},j,t^0_{*,j}} \\
  M_j^{b_j} \coloneq \dec{K_j^{b_j},C_j^{b_j}} \\
  \pcendfor{j} \\
  \pcreturn
  \pcnote{
    $M^0\coloneq\bigl(M_j^0\mid j\in[0,m)\bigr)$、
    $M^1\coloneq\bigl(M_j^1\mid j\in[0,m)\bigr)$ 为两组输入消息，
    $M^b\coloneq\bigl(M_j^{b_j}\mid j\in[0,m)\bigr)$ 为所选消息族。}
\end{protocol}

把 BaseOTSetup、IknpExpand、PayloadOT 编排成为完整的 IKNP03 协议。

\begin{protocol}[
  name=IknpOT,
  params={\party{S}: (M^0,M^1),~ \party{R}: b \mid \mathtt{sid}, \mathtt{tag},m,\kappa},
  returns={\party{S}: \varnothing,~ \party{R}: M^b},
  layout=multiline
]
  \pcparty{S,R}
  \party{S}: (s,k^s),~ \party{R}: (k^0,k^1) \coloneq
  \call{BaseOTSetup}{\kappa} \\
  \party{S}: (q,u),~ \party{R}: (t^0,u) \coloneq \\
  \quad \mathrm{IknpExpand}
  \left(\party{S}: (s,k^s),~ \party{R}: (k^0,k^1,b) \mid \mathtt{sid}, \mathtt{tag},m\right) \\
  \party{S}: \varnothing,~ \party{R}: M^b \coloneq \\
  \quad \mathrm{PayloadOT}
  \left(\party{S}: (q,s,M^0,M^1),~ \party{R}: (t^0,b) \mid \mathtt{sid}, \mathtt{tag},m\right) \\
  \pcreturn
\end{protocol}

\subsection{KOS15}

KOS 在发送 Payload 密文之前，检查矩阵各行是否使用一致的选择向量。
$\sum$ 表示域 $\mathbb{F}$ 元素的加法，即按位异或。
点乘 “$\cdot$” 表示域 $\mathbb{F}=\mathrm{GF}(2^\kappa)$ 元素的乘法，
现在（2026年9月）的 AI 已经能描述这种乘法如何实施，这里不再赘述。
单比特 $b_j$ 视为域中的常数 $0$ 或 $1$，故 $b_j\cdot\chi_j$ 表示选取或丢弃 $\chi_j$。

\begin{protocol}[
  name=KosCheck,
  params={\party{S}: (q,s),~ \party{R}: (t^0,b) \mid \chi,L},
  returns={\party{S, R}: \varnothing},
  layout=multiline
]
  \pcparty{R}
  \tau \coloneq \sum_{j=0}^{L-1}t^0_{*,j}\cdot\chi_j \in \mathbb F \\
  \rho \coloneq \sum_{j=0}^{L-1}b_j\cdot\chi_j \in \mathbb F \\
  \pcsend{\tau,\rho}{\party{S}} \\
  \pcparty{S}
  Q \coloneq \sum_{j=0}^{L-1}q_{*,j}\cdot\chi_j \in \mathbb F \\
  \pcif Q \ne \tau\oplus(\rho\cdot s); \pcthen
  \pccomment{请读者自行验算等式两边，建议交给 AI。}\\
  \pcabort \\
  \pcfi \\
  \pcreturn
\end{protocol}

\begin{protocol}[
  name=KosOT,
  params={\party{S}: (s,k^s,M^0,M^1),~ \party{R}: (k^0,k^1,b) \mid \mathtt{sid}, \mathtt{tag},m,m_2},
  returns={\party{S}: \varnothing,~ \party{R}: M^b},
  layout=multiline
]
  \pcparty{S,R}m'\coloneq m+\kappa+m_2 \\
  \pcparty{R}
  \tilde b\coloneq\uniform{\bits{\kappa+m_2}},~
  b'\coloneq b\mathbin{\|}\tilde b \\
  \pcparty{S,R}
  \party{S}: (q,u),~ \party{R}: (t^0,u) \coloneq \\
  \quad \call{IknpExpand}
  {\party{S}: (s,k^s),~ \party{R}: (k^0,k^1,b') \mid \mathtt{sid}, \mathtt{tag},m'} \\
  \chi\coloneq\prf{
    \texttt{"kos-check"},\mathtt{sid}, \mathtt{tag},\kappa,m,m_2,u
  }\in\mathbb F^{m'} \\
  \call{KosCheck}{\party{S}: (q,s),~ \party{R}: (t^0,b') \mid \chi,m'} \\
  \party{S}: \varnothing,~ \party{R}: M^b \coloneq \\
  \quad \call{PayloadOT}{
    \party{S}: (q,s,M^0,M^1),~ \party{R}: (t^0,b)
    \mid \mathtt{sid}, \mathtt{tag},m
  } \\
  \pcreturn
\end{protocol}

\paragraph{陪跑比特} KosOT 给实际选择向量 $b$ 追加 $\kappa+m_2$ 个随机陪跑比特，
目的是隐藏检查响应 $\rho$ 中关于 $b$ 的信息，$m_2$ 是统计安全参数。
检查覆盖全部 $m'=m+\kappa+m_2$ 列，PayloadOT 只使用前 $m$ 列；
陪跑列没有对应消息，也无需派生消息密钥。
这一安排及其安全分析直接见
\href{https://www.iacr.org/archive/crypto2015/92160327/92160327.pdf\#page=11}
{Keller、Orsini、Scholl，CRYPTO 2015，\S3.3 的 Lemma 2 及恶意 Sender 情形的证明}。
论文中的统计安全参数 $s$ 对应本文的 $m_2$。

\section{DKLS23 本身}

以下固定素数阶群 $\langle G\rangle$，群阶为 $n$，标量运算均在 $\mathbb Z_n$ 中进行。
$\kappa=\lceil\log_2n\rceil$，$\lambda$ 是种子长度，$\lambda_c,\lambda_s$ 分别是计算与统计安全参数。
默认参数为
\[
  \kappa=\lambda=256,\qquad \lambda_c=\lambda_s=128,
  \qquad k=4,\quad q=2^k=16.
\]

$H$ 表示带用途标签的随机预言机接口，输入先作含类型、长度和顺序的无歧义编码。
未注明时输出为 $\bits{2\lambda_c}$；注明域、矩阵形状或比特长度时，按该类型扩展输出，
并把目标类型纳入域分离。到 $\mathbb Z_n$ 的采样使用 XOF 产生候选标量并拒绝大于等于 $n$ 的值。
种子扩展记作 $\prf{d,K}\in\bits{\ell}$，其中 $K$ 为秘密种子，$d$ 为公开域分离输入。
所有 $\mathtt{sid}$ 均绑定协议版本、参数、双方身份、角色与唯一实例号。
并列随机赋值表示各变量独立采样；数组按索引递增排列，无下标的数组名表示相应的完整有序族。
接收消息时先验证上下文、长度、标量和群元素的规范编码及群成员资格。
$\call{Check}{P}$ 表示检查命题 $P$，不成立则立即中止；初始化与签名的失败处理范围在编排节规定。

\subsection{Endemic OT}

采用 \href{https://eprint.iacr.org/2019/706.pdf}{Masny--Rindal，Fig.\ 8}
所用的 Endemic OT 消息与密钥公式。
将本地计算拆开，便于在 Keygen 中与其他消息合并发送。
单独运行时，Receiver 的 $R^0,R^1$ 与 Sender 的 $V^0,V^1$ 可以同轮发送。

\begin{protocol}[
  name=EndemicRecvInit,
  params={s\in\bits{\kappa}\mid\mathtt{sid}},
  returns={u,R^0,R^1}
]
  \pcfor a\in[0,\kappa);\pcdo \\
  u_a\coloneq\uniform{\mathbb Z_n},~w_a\coloneq\uniform{\bits{\lambda}} \\
  R_a^{1-s_a}\coloneq\mathrm{HashToGroup}(\mathtt{sid},\texttt{"ot-r"},a,w_a) \\
  R_a^{s_a}\coloneq u_aG-\mathrm{HashToGroup}(\mathtt{sid},\texttt{"ot-u"},s_a,a,R_a^{1-s_a}) \\
  \pcendfor{a} \\
  \pcreturn
\end{protocol}

\begin{protocol}[
  name=EndemicSendInit,
  params={\kappa},
  returns={v^0,v^1,V^0,V^1}
]
  \pcfor a\in[0,\kappa),~d\in\bit;\pcdo \\
  v_a^d\coloneq\uniform{\mathbb Z_n},~V_a^d\coloneq v_a^dG \\
  \pcendfor{a,d} \\
  \pcreturn
\end{protocol}

\begin{protocol}[
  name=EndemicSendFinish,
  params={v^0,v^1,R^0,R^1\mid\mathtt{sid}},
  returns={\rho^0,\rho^1}
]
  \pcfor a\in[0,\kappa),~d\in\bit;\pcdo \\
  U_a^d\coloneq R_a^d+\mathrm{HashToGroup}(\mathtt{sid},\texttt{"ot-u"},d,a,R_a^{1-d}) \\
  \rho_a^d\coloneq H(\mathtt{sid},\texttt{"ot-key"},d,a,v_a^dU_a^d)\in\bits{\lambda} \\
  \pcendfor{a,d} \\
  \pcreturn
\end{protocol}

\begin{protocol}[
  name=EndemicRecvFinish,
  params={s,u,V^0,V^1\mid\mathtt{sid}},
  returns={\rho^s}
]
  \pcfor a\in[0,\kappa);\pcdo \\
  \rho_a^{s_a}\coloneq H(\mathtt{sid},\texttt{"ot-key"},s_a,a,u_aV_a^{s_a})\in\bits{\lambda} \\
  \pcendfor{a} \\
  \pcreturn
\end{protocol}

这里 $\mathrm{HashToGroup}$ 使用离散对数未知的点哈希，不能实现为 $H(\cdots)G$；
$w_a$ 新鲜且保密。正确性来自 $U_a^{s_a}=u_aG$。
Endemic OT 允许恶意方影响自己所得密钥的分布，后续组合须保留这种接口条件，
不能把它直接当作双方输出均不可操纵的随机 OT。

\subsection{PPRF}

采用 \href{https://eprint.iacr.org/2022/192.pdf}{Roy，SoftSpokenOT，Fig.\ 13--14}
的建树、打孔与叶子校验步骤。每棵树消耗 $k$ 个 Base OT；
完整叶子族记为 $A$，缺少位置 $\delta$ 的叶子族记为 $A^*$。
$\mathtt{tid}$ 绑定初始化会话及树编号；本节 PRF 的域分离输入均包含 $\mathtt{tid}$。

\begin{protocol}[
  name=PprfBuild,
  params={\rho^0,\rho^1\mid\mathtt{tid},k},
  returns={A,P},
  layout=multiline
]
  T^1_0\coloneq\rho_0^0,~T^1_1\coloneq\rho_0^1 \\
  \pcfor a\in[1,k);\pcdo \\
  \pcfor z\in[0,2^a);\pcdo \\
  T^{a+1}_{2z}\mathbin{\|}T^{a+1}_{2z+1}
    \coloneq\prf{\mathtt{tid},\texttt{"tree"},T^a_z}\in\bits{2\lambda} \\
  \pcendfor{z} \\
  t_a^d\coloneq\rho_a^d\oplus\bigoplus_{z=0}^{2^a-1}T^{a+1}_{2z+d},\quad d\in\bit \\
  \pcendfor{a} \\
  A\coloneq(T^k_z\mid z\in[0,q)) \\
  \tilde s_z\coloneq\prf{\mathtt{tid},\texttt{"leaf"},T^k_z}\in\bits{2\lambda},\quad z\in[0,q) \\
  \tilde t\coloneq\bigoplus_{z=0}^{q-1}\tilde s_z \\
  \tilde s\coloneq H(\mathtt{tid},\texttt{"leaves"},(\tilde s_z\mid z\in[0,q)))\in\bits{2\lambda} \\
  P\coloneq(t^0,t^1,\tilde t,\tilde s) \\
  \pcreturn
\end{protocol}

\begin{protocol}[
  name=PprfEval,
  params={s,\rho^s,P\mid\mathtt{tid},k},
  returns={\delta,A^*}
]
  (t^0,t^1,\tilde t,\tilde s)\coloneq P \\
  y_1\coloneq1-s_0,~T^1_{s_0}\coloneq\rho_0^{s_0} \\
  \pcfor a\in[1,k);\pcdo \\
  \pcfor z\in[0,2^a)\setminus\{y_a\};\pcdo \\
  T^{a+1}_{2z}\mathbin{\|}T^{a+1}_{2z+1}
    \coloneq\prf{\mathtt{tid},\texttt{"tree"},T^a_z}\in\bits{2\lambda} \\
  \pcendfor{z} \\
  T^{a+1}_{2y_a+s_a}\coloneq t_a^{s_a}\oplus\rho_a^{s_a}
    \oplus\bigoplus_{z\ne y_a}T^{a+1}_{2z+s_a} \\
  y_{a+1}\coloneq2y_a+1-s_a \\
  \pcendfor{a} \\
  \delta\coloneq y_k,~A^*\coloneq(T^k_z\mid z\in[0,q)\setminus\{\delta\}) \\
  \pcfor z\in[0,q)\setminus\{\delta\};\pcdo \\
  \tilde s_z\coloneq\prf{\mathtt{tid},\texttt{"leaf"},T^k_z}\in\bits{2\lambda} \\
  \pcendfor{z} \\
  \tilde s_\delta\coloneq\tilde t\oplus\bigoplus_{z\ne\delta}\tilde s_z \\
  \call{Check}{H(\mathtt{tid},\texttt{"leaves"},(\tilde s_z\mid z\in[0,q)))=\tilde s} \\
  \pcreturn
\end{protocol}

Sender 发送 $P$，Receiver 调用 PprfEval；校验通过后才保存初始化结果。
Base OT 选择位与打孔位置的关系为
\[
  \delta=\sum_{a=0}^{k-1}(1-s_a)2^{k-1-a}.
\]
后文统一按高位在前编码 $\delta$，即 $\operatorname{bit}_a(\delta)=1-s_a$。
每个有向初始化还保存 Base OT 与 PPRF 公共消息的规范 transcript 摘要 $T_{\mathrm{init}}$，
包含会话、身份、角色、参数及版本；消息按协议规定的顺序编码，不按本地到达顺序编码。

\subsection{SoftSpokenOT}

完整树持有者成为扩展 OT 的 Receiver，打孔树持有者成为扩展 OT 的 Sender。
每批使用 $h\in[0,\kappa/k)$ 的树，扩展出 $L$ 个随机 OT。
以下将 \href{https://eprint.iacr.org/2022/192.pdf}{SoftSpokenOT}
的展开与一致性检查合为一条 Receiver 消息，使用 05 中的随机预言机挑战实现。

令 $\mathbb F_c=\mathrm{GF}(2^\sigma)$，固定其多项式基，默认 $\sigma=128$。
$L$ 是 $\sigma$ 的倍数，$B=L/\sigma$，$L'=L+\sigma$。
将 $x\in\bits{L'}$ 分为 $B+1$ 个域元素 $x^{[0]},\ldots,x^{[B]}$，定义
\[
  h_\chi(x)\coloneq x^{[B]}+\sum_{b=0}^{B-1}\chi_bx^{[b]}
  \in\mathbb F_c.
\]
本节域加法为异或；$\odot$ 表示由单比特控制的取值或置零。

\begin{protocol}[
  name=SoftSpokenRecv,
  params={A,\boldsymbol\beta\in\bits{L}\mid\mathtt{sid},L},
  returns={\rho^{\boldsymbol\beta},E},
  layout=multiline
]
  \beta^{\mathrm{ext}}\coloneq\uniform{\bits{\sigma}},~
    \hat\beta\coloneq\boldsymbol\beta\mathbin{\|}\beta^{\mathrm{ext}} \\
  \pcfor h\in[0,\kappa/k);\pcdo \\
  r_{h,x}\coloneq\prf{\mathtt{sid},h,\texttt{"expand"},A_{h,x}}\in\bits{L'},\quad x\in[0,q) \\
  u_h\coloneq\hat\beta\oplus\bigoplus_{x=0}^{q-1}r_{h,x} \\
  \pcfor a\in[0,k);\pcdo \\
  v_{hk+a,*}\coloneq\bigoplus_{x=0}^{q-1}\operatorname{bit}_a(x)\odot r_{h,x} \\
  \pcendfor{a} \\
  \pcendfor{h} \\
  \chi\coloneq H(\mathtt{sid},\texttt{"ss-check"},u)\in\mathbb F_c^B \\
  \tau_a\coloneq h_\chi(v_{a,*}),\quad a\in[0,\kappa) \\
  \tilde\beta\coloneq h_\chi(\hat\beta),~E\coloneq(u,\tau,\tilde\beta) \\
  \rho_j^{\beta_j}\coloneq H(\mathtt{sid},\texttt{"ss-key"},j,v_{*,j})\in\bits{\lambda},\quad j\in[0,L) \\
  \pcreturn
\end{protocol}

\begin{protocol}[
  name=SoftSpokenSend,
  params={\delta,A^*,E\mid\mathtt{sid},L},
  returns={\rho^0,\rho^1}
]
  (u,\tau,\tilde\beta)\coloneq E \\
  \chi\coloneq H(\mathtt{sid},\texttt{"ss-check"},u)\in\mathbb F_c^B \\
  \pcfor h\in[0,\kappa/k);\pcdo \\
  \pcfor x\in[0,q)\setminus\{\delta_h\};\pcdo \\
  r_{h,x}\coloneq\prf{\mathtt{sid},h,\texttt{"expand"},A^*_{h,x}}\in\bits{L'} \\
  \pcendfor{x} \\
  \pcfor a\in[0,k);\pcdo \\
  \Delta_{hk+a}\coloneq\operatorname{bit}_a(\delta_h) \\
  w_{hk+a,*}\coloneq\bigoplus_{x\ne\delta_h}\operatorname{bit}_a(\delta_h\oplus x)\odot r_{h,x} \\
  w_{hk+a,*}\coloneq w_{hk+a,*}\oplus(\Delta_{hk+a}\odot u_h) \\
  \call{Check}{h_\chi(w_{hk+a,*})=\tau_{hk+a}\oplus(\Delta_{hk+a}\odot\tilde\beta)} \\
  \pcendfor{a} \\
  \pcendfor{h} \\
  \pcfor j\in[0,L);\pcdo \\
  \rho_j^0\coloneq H(\mathtt{sid},\texttt{"ss-key"},j,w_{*,j})\in\bits{\lambda} \\
  \rho_j^1\coloneq H(\mathtt{sid},\texttt{"ss-key"},j,w_{*,j}\oplus\Delta)\in\bits{\lambda} \\
  \pcendfor{j} \\
  \pcreturn
\end{protocol}

Receiver 发送 $E$，Sender 验证后才派生输出密钥。
诚实执行满足 $w_{*,j}=v_{*,j}\oplus(\hat\beta_j\odot\Delta)$；
最后 $\sigma$ 列只参与检查，输出仅使用前 $L$ 列。
每批保存
\[
  T_{\mathrm{OT}}\coloneq H(\mathtt{sid},\texttt{"ot-transcript"},T_{\mathrm{init}},E).
\]
初始化摘要不能由当前扩展消息代替。
挑战的具体 transcript 绑定沿用
\href{https://eprint.iacr.org/2026/976.pdf}{2026/976，\S4.5} 的组合要求；
Endemic OT、PPRF、扩展 OT 的组合还须满足该处的输入绑定条件。

\subsection{RVOLE}

保留 \href{https://eprint.iacr.org/2026/976.pdf}{2026/976，\S4}
的完整一致性检查和随机输入接口，指定输入通过偏移接入。
功能负载数 $N_1=2$，默认校验负载数 $N_2=1$；
协议暂按一般的 $N_1,N_2$ 书写，令 $N=N_1+N_2$。
Receiver 为每个实例独立采样 $\boldsymbol\beta\in\bits{L}$，
经 SoftSpokenOT 获得 $\rho^{\boldsymbol\beta}$，Sender 获得 $\rho^0,\rho^1$。
二者使用同一个 $T_{\mathrm{OT}}$。

定义域上的密钥扩展
\[
  \alpha^d_{j,a}\coloneq\mathrm{PRF}_{\rho_j^d}
  \circ H(\mathtt{sid},\texttt{"ot-pad"},j,a)\in\mathbb Z_n.
\]
这里的 PRF 输出也按前述拒绝采样规则产生标量。
RvoleSend 先生成消息及逐行份额，RvoleRecv 验证后产生逐行份额；
最后的 gadget 聚合单独列出，以兼容两种生成时机。

\begin{protocol}[
  name=RvoleSend,
  params={\rho^0,\rho^1\mid\mathtt{sid},T_{\mathrm{OT}},L,N_1,N_2},
  returns={w,\alpha^0,Q},
  layout=multiline
]
  w_a\coloneq\uniform{\mathbb Z_n},\quad a\in[0,N) \\
  \pcfor j\in[0,L),~a\in[0,N);\pcdo \\
  \alpha^d_{j,a}\coloneq\mathrm{PRF}_{\rho_j^d}\circ
    H(\mathtt{sid},\texttt{"ot-pad"},j,a)\in\mathbb Z_n,\quad d\in\bit \\
  \tilde a_{j,a}\coloneq\alpha^0_{j,a}-\alpha^1_{j,a}+w_a \\
  \pcendfor{j,a} \\
  \Theta\coloneq H(\mathtt{sid},\texttt{"challenge"},T_{\mathrm{OT}},\tilde a)
    \in\mathbb Z_n^{N_2\times N_1} \\
  \pcfor a\in[N_1,N);\pcdo \\
  \eta_a\coloneq w_a+\sum_{b=0}^{N_1-1}\theta_{a,b}w_b \\
  \mu_{j,a}\coloneq\alpha^0_{j,a}+\sum_{b=0}^{N_1-1}\theta_{a,b}\alpha^0_{j,b},\quad j\in[0,L) \\
  \pcendfor{a} \\
  \mu\coloneq H(\mathtt{sid},\texttt{"check"},(\mu_{j,a}\mid j\in[0,L),~a\in[N_1,N))) \\
  Q\coloneq(\tilde a,\eta,\mu) \\
  w\coloneq(w_b\mid b\in[0,N_1)) \\
  \pcreturn
\end{protocol}

\begin{protocol}[
  name=RvoleRecv,
  params={\boldsymbol\beta,\rho^{\boldsymbol\beta},Q\mid\mathtt{sid},T_{\mathrm{OT}},L,N_1,N_2},
  returns={t},
  layout=multiline
]
  (\tilde a,\eta,\mu)\coloneq Q \\
  \pcfor j\in[0,L),~a\in[0,N);\pcdo \\
  \alpha^{\beta_j}_{j,a}\coloneq\mathrm{PRF}_{\rho_j^{\beta_j}}\circ
    H(\mathtt{sid},\texttt{"ot-pad"},j,a)\in\mathbb Z_n \\
  t_{j,a}\coloneq\alpha^{\beta_j}_{j,a}+\beta_j\tilde a_{j,a} \\
  \pcendfor{j,a} \\
  \Theta\coloneq H(\mathtt{sid},\texttt{"challenge"},T_{\mathrm{OT}},\tilde a)
    \in\mathbb Z_n^{N_2\times N_1} \\
  \pcfor j\in[0,L),~a\in[N_1,N);\pcdo \\
  \hat\mu_{j,a}\coloneq t_{j,a}+\sum_{b=0}^{N_1-1}\theta_{a,b}t_{j,b}-\beta_j\eta_a \\
  \pcendfor{j,a} \\
  \hat\mu\coloneq H(\mathtt{sid},\texttt{"check"},(\hat\mu_{j,a}\mid j\in[0,L),~a\in[N_1,N))) \\
  \call{Check}{\hat\mu=\mu} \\
  \pcreturn
\end{protocol}

矩阵与族均按索引递增编码；$\Theta$ 的行按 $a\in[N_1,N)$ 编号。
检查通过后再使用功能负载，校验负载不输出给签名层。
诚实执行时 $t_{j,a}=\alpha^0_{j,a}+\beta_jw_a$，因此检查中所有 $\beta_j$ 项抵消。
这条正确性等式不替代恶意安全分析，后者还要求相应的随机 OT 或具体 transcript 输入绑定性质。

\subsection{gadget 合成与分解}

gadget 将较长的选择串合成为域上的标量，同时聚合双方的逐行乘法份额。
采用 \href{https://eprint.iacr.org/2026/976.pdf}{2026/976，\S4--5}
中适用于两种编排的方案：多方流程取 III，两方独立预处理取 II。
二者使用同一组 RVOLE 消息公式，只改变 $g$ 的产生时机与 $L$。

\begin{protocol}[
  name=Gadget,
  params={v,Q\mid\mathtt{sid},T_{\mathrm{OT}},L},
  returns={g\in\mathbb Z_n^L}
]
  \pcif v=\mathrm{III};\pcthen \\
  g\coloneq H(\mathtt{sid},\texttt{"gadget-III"})\in\mathbb Z_n^L \\
  \pcelse \\
  \call{Check}{v=\mathrm{II}\ \land\ Q\ne\bot} \\
  D_Q\coloneq H(\mathtt{sid},\texttt{"rvole-msg"},Q) \\
  g\coloneq H(\mathtt{sid},\texttt{"gadget-II"},T_{\mathrm{OT}},D_Q)\in\mathbb Z_n^L \\
  \pcfi \\
  \pcreturn
\end{protocol}

III 在实例开始时调用，允许输入 $Q=\bot$。
II 必须等 RvoleSend 的最终消息 $Q$ 固定后调用；Receiver 先验证 $Q$，再输出聚合结果。
III 的 $T_{\mathrm{OT}}$ 尚未产生时也可传入 $\bot$，该分支不使用它。

\begin{protocol}[
  name=RvoleAggregate,
  params={\party{S}: (w,\alpha^0),~\party{R}: (\boldsymbol\beta,t)\mid g,L,N_1},
  returns={\party{S}: (w,y),~\party{R}: (\beta,z)},
  st={\forall a\in[0,N_1),~y_a+z_a=w_a\beta},
  layout=multiline
]
  \pcparty{R}\beta\coloneq\sum_{j=0}^{L-1}g_j\beta_j \\
  \pcparty{S,R}\pcfor a\in[0,N_1);\pcdo \\
  \pcparty{S}y_a\coloneq-\sum_{j=0}^{L-1}g_j\alpha^0_{j,a} \\
  \pcparty{R}z_a\coloneq\sum_{j=0}^{L-1}g_jt_{j,a} \\
  \pcendfor{a} \\
  \pcreturn
\end{protocol}

指定 Sender 输入 $a$ 时，使用以下偏移；$w$ 必须是本次新鲜采样的随机量。
Sender 若已经知道 $a$，可将偏移与 $Q$ 同轮发送。

\begin{protocol}[
  name=RvoleShift,
  params={\party{S}: (a,w,y),~\party{R}: (\beta,z)},
  returns={\party{S}: y,~\party{R}: z},
  st={y+z=a\beta},
  layout=multiline
]
  \pcparty{S}\delta\coloneq a-w \\
  \pcsend{\delta}{\party{R}} \\
  \pcparty{R}z\coloneq z+\beta\delta \\
  \pcreturn
\end{protocol}

所谓分解，是在生成时保留 $\boldsymbol\beta$ 以及关系 $\beta=\langle g,\boldsymbol\beta\rangle$；
此处没有把任意给定标量高效逆解成选择串的步骤。
对于指定的 Sender 输入，用 RvoleShift 接入即可。

按上述文献的充分条件，II 要求 $L\ge\kappa+2\lambda_c+3\lambda_s$；III 要求
\[
  L\ge\kappa+7\lambda_s+2\log_2(\lambda_s+1)
     +2\lambda_s\log_2(eL/\lambda_s).
\]
取单个校验负载还要求 $\log_2n\ge\lambda_c+2\log_2L$。
对前述默认参数，按 $\sigma=128$ 对齐后使用
\[
  \begin{array}{c|rrrr}
    \text{用途}&L&L'&B&N_2\\\hline
    \text{多方，III}&2688&2816&21&1\\
    \text{两方独立预处理，II}&896&1024&7&1
  \end{array}
\]
两种方案都使用 $256$ 个 Base OT、$64$ 棵深度为 $4$ 的树。
这些是相应 RVOLE 变体的参数条件，整套签名的安全结论还取决于底层实现与组合前提。

\subsection{编排}

参与方集合 $\mathbb K\subset\mathbb Z_n^*$，各编号互异且非零；
阈值为 $t$，签名集合 $\mathbb S\subseteq\mathbb K$ 满足 $|\mathbb S|\ge t$。
标量与秘密种子通过认证保密信道传输，广播消息须有一致的发送者身份与会话标识。
下列伪代码在每方 $i$ 执行，$j$ 始终表示其他参与方；
每轮收齐所需消息并通过检查后，才进入下一轮。

树的下标统一按 PPRF Sender、Receiver 排列：$A_{h,i,j}$ 由 $i$ 持有，
对应的 $(\delta_{h,i,j},A^*_{h,i,j})$ 由 $j$ 持有。
扩展 OT 和 RVOLE 则按各自的 Sender、Receiver 排列下标。
因此，有向 RVOLE $i\to j$ 使用 $j$ 的完整树 $A_{h,j,i}$ 和 $i$ 的打孔树 $A^*_{h,j,i}$，
其 SoftSpoken 消息 $E_{i,j}$ 由 $j$ 发给 $i$，RVOLE 消息 $Q_{i,j}$ 由 $i$ 发给 $j$。

\paragraph{失败与状态} Keygen 中任何检查失败均中止初始化，销毁临时秘密，不输出 Keyshare。
签名中任何失败均废弃本次预签名，重试使用新会话、新随机数和新选择串。
RVOLE、逐对 $\Gamma$、SoftSpoken 或 nonce 承诺检查失败时，还须持久停用相应对端，
取消所有共享该 OT 状态且涉及该对端的并发会话与待用预签名。
这个停用记录须在共享状态的进程和设备间生效，更换会话号不能清除它。

下文用 $\call{CheckP}{P,j}$ 表示检查失败时停用对端 $j$，其余 Check 只中止当前执行。
对无法定位的承诺视图或聚合公钥不一致，沿用 07 的保守策略，停用 $\mathbb S\setminus\{i\}$；
这是停止交互的范围，不是公开归责结论。
收到失败通知或因超时取消时，只终止相应执行，不仅凭通知或超时停用某方。
失败通知绑定会话且不含秘密；成功、失败、取消后的临时材料都按相同生命周期销毁。
RVOLE 与 $\Gamma$ 失败后的停止交互要求见
\href{https://eprint.iacr.org/2023/765.pdf}{DKLS23，Protocol 3.6，第 8--10 步}；
共享 OT 状态的限制见 \href{https://eprint.iacr.org/2026/976.pdf}{2026/976，\S2.1}。

\subsubsection{Keygen}

这里保留 07 所用的 Feldman 承诺与 DLog 知识证明版本。
它的实现来源尚未逐项对应论文，不能套用
\href{https://eprint.iacr.org/2023/765.pdf}{DKLS23，Protocol 7.1}
对 relaxed Keygen 的证明。以下将其实际计算、校验及 OT 初始化写全。

对公开点 $X=xG$ 的知识证明采用 Schnorr 形式。
证明上下文 $c$ 绑定 Keygen 会话、证明者及系数编号或最终份额用途。
证明者取 $a\gets\uniform{\mathbb Z_n}$，令
\[
  A=aG,\qquad e=H(c,\texttt{"dlog"},X,A)\in\mathbb Z_n,
  \qquad \pi=(A,z),\quad z=a+ex.
\]
$\mathrm{DLogProve}(x\mid c)$ 返回 $\pi$；
$\mathrm{DLogVerify}(X,\pi\mid c)$ 在验证编码后，重算 $e$ 并返回 $zG=A+eX$ 的真假。
多项式点承诺的求值和 Lagrange 系数记为
\[
  F_i(z)\coloneq\sum_{a=0}^{t-1}z^aF_{i,a},\qquad
  \lambda(i,J)\coloneq\prod_{j\in J\setminus\{i\}}\frac{-j}{i-j}.
\]

\begin{protocol}[
  name=Keygen,
  params={i\in\mathbb K\mid\mathtt{kgid},\mathbb K,t},
  returns={\mathrm{Keyshare}_i},
  layout=multiline
]
  c_{i,a}\coloneq\uniform{\mathbb Z_n},~F_{i,a}\coloneq c_{i,a}G,\quad a\in[0,t) \\
  P_i(X)\coloneq\sum_{a=0}^{t-1}c_{i,a}X^a \\
  e_i\coloneq\uniform{\bits{\lambda}},~C_i\coloneq H(\mathtt{kgid},\texttt{"feldman"},i,F_i,e_i) \\
  \pcfor j\in\mathbb K\setminus\{i\};\pcdo \\
  \epsilon^\circ_{i,j},\nu_{i,j}\coloneq\uniform{\bits{\lambda}} \\
  C^\epsilon_{i,j}\coloneq H(\mathtt{kgid},\texttt{"zero-seed"},i,j,\epsilon^\circ_{i,j},\nu_{i,j}) \\
  \pcsend{C_i,C^\epsilon_{i,j}}{j}\pccomment{第一轮} \\
  \pcendfor{j} \\
  \pi_{i,a}\coloneq\mathrm{DLogProve}(c_{i,a}\mid\mathtt{kgid},i,\texttt{"coefficient"},a),\quad a\in[0,t) \\
  \pcfor j\in\mathbb K\setminus\{i\};\pcdo \\
  b_{j,i}\coloneq H(\mathtt{kgid},\texttt{"base"},j,i) \\
  s_{j,i}\coloneq\uniform{\bits{\kappa}} \\
  (u_{j,i},R^0_{j,i},R^1_{j,i})\coloneq\call{EndemicRecvInit}{s_{j,i}\mid b_{j,i}} \\
  \pcsend{F_i,e_i,\pi_i,R^0_{j,i},R^1_{j,i},\epsilon^\circ_{i,j},\nu_{i,j}}{j}
    \pccomment{第二轮} \\
  \pcendfor{j} \\
  \pcfor j\in\mathbb K\setminus\{i\};\pcdo \\
  \call{Check}{C_j=H(\mathtt{kgid},\texttt{"feldman"},j,F_j,e_j)} \\
  \call{Check}{C^\epsilon_{j,i}=H(\mathtt{kgid},\texttt{"zero-seed"},j,i,\epsilon^\circ_{j,i},\nu_{j,i})} \\
  \call{Check}{\forall a\in[0,t),~\mathrm{DLogVerify}(F_{j,a},\pi_{j,a}\mid\mathtt{kgid},j,\texttt{"coefficient"},a)} \\
  \epsilon_{i,j}\coloneq\epsilon^\circ_{i,j}\oplus\epsilon^\circ_{j,i} \\
  \pcendfor{j} \\
  F\coloneq\sum_{j\in\mathbb K}F_j,~d_{i,j}\coloneq P_i(j),\quad j\in\mathbb K \\
  \pcfor j\in\mathbb K\setminus\{i\};\pcdo \\
  b_{i,j}\coloneq H(\mathtt{kgid},\texttt{"base"},i,j) \\
  (v^0,v^1,V^0_{i,j},V^1_{i,j})\coloneq\call{EndemicSendInit}{\kappa} \\
  (\rho^0_{i,j},\rho^1_{i,j})\coloneq\call{EndemicSendFinish}{v^0,v^1,R^0_{i,j},R^1_{i,j}\mid b_{i,j}} \\
  \pcfor h\in[0,\kappa/k);\pcdo \\
  \mathtt{tid}_{h,i,j}\coloneq H(b_{i,j},\texttt{"pprf"},h) \\
  (A_{h,i,j},P_{h,i,j})\coloneq \\
  \quad\call{PprfBuild}{\rho^0_{i,j}[hk:(h+1)k],\rho^1_{i,j}[hk:(h+1)k]\mid\mathtt{tid}_{h,i,j},k} \\
  \pcendfor{h} \\
  \pcsend{d_{i,j},F,V^0_{i,j},V^1_{i,j},(P_{h,i,j}\mid h\in[0,\kappa/k))}{j}
    \pccomment{第三轮} \\
  \pcendfor{j} \\
  \pcfor j\in\mathbb K\setminus\{i\};\pcdo \\
  \call{Check}{F^{(j)}=F\ \land\ d_{j,i}G=F_j(i)} \\
  \rho^{s_{j,i}}_{j,i}\coloneq\call{EndemicRecvFinish}{s_{j,i},u_{j,i},V^0_{j,i},V^1_{j,i}\mid b_{j,i}} \\
  \pcfor h\in[0,\kappa/k);\pcdo \\
  \mathtt{tid}_{h,j,i}\coloneq H(b_{j,i},\texttt{"pprf"},h) \\
  (s',\rho')\coloneq(s_{j,i}[hk:(h+1)k],\rho^{s_{j,i}}_{j,i}[hk:(h+1)k]) \\
  (\delta_{h,j,i},A^*_{h,j,i})\coloneq\call{PprfEval}{s',\rho',P_{h,j,i}\mid\mathtt{tid}_{h,j,i},k} \\
  \pcendfor{h} \\
  T^{\mathrm{init}}_{i,j}\coloneq H(b_{i,j},\texttt{"init"},R^0_{i,j},R^1_{i,j},V^0_{i,j},V^1_{i,j},P_{*,i,j}) \\
  T^{\mathrm{init}}_{j,i}\coloneq H(b_{j,i},\texttt{"init"},R^0_{j,i},R^1_{j,i},V^0_{j,i},V^1_{j,i},P_{*,j,i}) \\
  \pcendfor{j} \\
  x_i\coloneq\sum_{j\in\mathbb K}d_{j,i},~Y\coloneq F(0),~S_i\coloneq x_iG \\
  \pi_i\coloneq\mathrm{DLogProve}(x_i\mid\mathtt{kgid},i,\texttt{"share"}) \\
  \pcsend{S_i,\pi_i}{\mathbb K\setminus\{i\}}\pccomment{第四轮} \\
  \call{Check}{\forall j\ne i,~\mathrm{DLogVerify}(S_j,\pi_j\mid\mathtt{kgid},j,\texttt{"share"})} \\
  \call{Check}{Y\ne\mathcal O\ \land\ Y=\sum_{j\in\mathbb K}\lambda(j,\mathbb K)S_j} \\
  \mathrm{Keyshare}_i\coloneq(i,x_i,Y,\mathbb K,t,\mathtt{kgid},A_{*,i,*},\delta_{*,*,i},A^*_{*,*,i},T^{\mathrm{init}},\epsilon_{i,*}) \\
  \pcreturn
\end{protocol}

切片 $v[a:b]$ 表示保留 $a\le j<b$ 的分量，并从零重新编号供子协议使用。
$F^{(j)}$ 是 $j$ 发来的聚合承诺，$P_{*,i,j}$ 按树编号排序。
Keyshare 另外保存协议参数、密钥标识和初始为空的持久停用记录；
输出后销毁多项式系数、份额传输副本、Base OT 临时指数及不再需要的中间种子。

\subsubsection{多方 PreSign：前两轮}

第一轮之前固定消息哈希 $m$、派生偏移 $\nabla x$ 及参与集合 $\mathbb S$。
PreSign 只表示三轮签名的前两轮；其输出已经绑定消息和公钥。
多方流程采用 Variant III，参数为 $L=2688,N_1=2,N_2=1$。
各方先检查参与集合、Keyshare 状态及停用记录，并固定
\[
  \begin{aligned}
    Y'&\coloneq Y+\nabla xG,\\
    C&\coloneq H(\texttt{"sign"},\mathtt{sid},\mathrm{keyid},\mathrm{sort}(\mathbb S),Y',m,\\
      &\hspace{4em}\mathrm{params},\mathrm{zeroMode},\mathrm{signMode}).
  \end{aligned}
\]
所有后续消息都携带 $C$，输入不符则中止。
有向子实例号不得重复，且必须绑定双方角色与调用编号。
$g_{i,j}$ 由 Gadget 的 III 分支在该实例开始时产生；收到扩展消息后计算 $T_{i,j}$。
相应公式为
\[
  \begin{aligned}
    c_{i,j}&\coloneq H(C,\texttt{"rvole"},i,j,\mathrm{callid}),\\
    T_{i,j}&\coloneq H(c_{i,j},\texttt{"ot-transcript"},T^{\mathrm{init}}_{j,i},E_{i,j}).
  \end{aligned}
\]

\begin{protocol}[
  name=PreSignRoundOne,
  params={\mathrm{Keyshare}_i,m,\nabla x,\mathbb S\mid C},
  returns={\mathrm{St}_i},
  layout=multiline
]
  r_i,\phi_i\coloneq\uniform{\mathbb Z_n^*},~R_i\coloneq r_iG \\
  \pcfor j\in\mathbb S\setminus\{i\};\pcdo \\
  I\coloneq\min(i,j),~J\coloneq\max(i,j) \\
  \zeta_{i,j}\coloneq H(C,\texttt{"zero-mask"},I,J,\epsilon_{i,j})\in\mathbb Z_n \\
  \pcendfor{j} \\
  \zeta_i\coloneq\sum_{j<i}\zeta_{i,j}-\sum_{j>i}\zeta_{i,j} \\
  x'_i\coloneq\lambda(i,\mathbb S)x_i+\zeta_i+\nabla x\,|\mathbb S|^{-1},~Y_i\coloneq x'_iG \\
  e_i^R\coloneq\uniform{\bits{\lambda}},~C_i^R\coloneq H(C,\texttt{"nonce"},i,R_i,e_i^R) \\
  \pcfor j\in\mathbb S\setminus\{i\};\pcdo \\
  \boldsymbol\beta_{j,i}\coloneq\uniform{\bits{L}},~\beta_{j,i}\coloneq\langle g_{j,i},\boldsymbol\beta_{j,i}\rangle \\
  (\rho^{\boldsymbol\beta_{j,i}}_{j,i},E_{j,i})\coloneq
    \call{SoftSpokenRecv}{A_{*,i,j},\boldsymbol\beta_{j,i}\mid c_{j,i},L} \\
  \pcsend{C_i^R,E_{j,i}}{j} \\
  \pcendfor{j} \\
  \mathrm{St}_i\coloneq\text{本轮本地变量、上下文及收发消息} \\
  \pcreturn
\end{protocol}

\begin{protocol}[
  name=PreSignRoundTwo,
  params={\mathrm{St}_i},
  returns={\mathrm{St}_i}
]
  d_i\coloneq H(C,\texttt{"nonce-view"},(C_j^R\mid j\in\mathrm{sort}(\mathbb S))) \\
  \pcfor j\in\mathbb S\setminus\{i\};\pcdo \\
  (\rho^0_{i,j},\rho^1_{i,j})\coloneq
    \call{SoftSpokenSend}{\delta_{*,j,i},A^*_{*,j,i},E_{i,j}\mid c_{i,j},L} \\
  (w_{i,j},\alpha^0_{i,j},Q_{i,j})\coloneq \\
  \quad\call{RvoleSend}{\rho^0_{i,j},\rho^1_{i,j}\mid c_{i,j},T_{i,j},L,2,1} \\
  y^a_{i,j}\coloneq-\sum_{\ell=0}^{L-1}g_{i,j,\ell}\alpha^0_{i,j,\ell,a},\quad a\in\{\mathrm r,\mathrm x\} \\
  \delta_{i,j}\coloneq(r_i,x'_i)-w_{i,j} \\
  \Gamma^{\mathrm r}_{i,j}\coloneq y^{\mathrm r}_{i,j}G,~\Gamma^{\mathrm x}_{i,j}\coloneq y^{\mathrm x}_{i,j}G \\
  \psi_{i,j}\coloneq\phi_i-\beta_{j,i} \\
  \pcsend{Q_{i,j},\delta_{i,j},R_i,e_i^R,Y_i,\Gamma^{\mathrm r}_{i,j},\Gamma^{\mathrm x}_{i,j},\psi_{i,j},d_i}{j} \\
  \pcendfor{j} \\
  \pcreturn
\end{protocol}

St 表示跨轮保留的本地状态，下一阶段继续使用并更新其中的变量；
上标 $\mathrm r,\mathrm x$ 分别对应 RVOLE 的功能负载 $0,1$。
SoftSpokenSend 失败时按 CheckP 处理其对端；RvoleSend 中的随机负载仅供本次使用。
接收第二轮后执行下列收尾，此处的 $\delta_{j,i}$ 是输入偏移，与树的打孔位置按索引形状区分。

\begin{protocol}[
  name=PreSignFinish,
  params={\mathrm{St}_i},
  returns={\mathrm{PreSignature}_i}
]
  \pcfor j\in\mathbb S\setminus\{i\};\pcdo \\
  \call{CheckP}{H(C,\texttt{"nonce"},j,R_j,e_j^R)=C_j^R,j} \\
  \call{CheckP}{d_j=d_i,\mathbb S\setminus\{i\}} \\
  t_{j,i}\coloneq\call{RvoleRecv}{\boldsymbol\beta_{j,i},\rho^{\boldsymbol\beta_{j,i}}_{j,i},Q_{j,i}\mid c_{j,i},T_{j,i},L,2,1} \\
  \pcfor a\in\{\mathrm r,\mathrm x\};\pcdo \\
  z^a_{j,i}\coloneq\sum_{\ell=0}^{L-1}g_{j,i,\ell}t_{j,i,\ell,a}+\beta_{j,i}\delta^a_{j,i} \\
  \pcendfor{a} \\
  \call{CheckP}{\beta_{j,i}R_j=\Gamma^{\mathrm r}_{j,i}+z^{\mathrm r}_{j,i}G,j} \\
  \call{CheckP}{\beta_{j,i}Y_j=\Gamma^{\mathrm x}_{j,i}+z^{\mathrm x}_{j,i}G,j} \\
  \pcendfor{j} \\
  \call{CheckP}{\sum_{j\in\mathbb S}Y_j=Y',\mathbb S\setminus\{i\}} \\
  R\coloneq\sum_{j\in\mathbb S}R_j,~\call{Check}{R\ne\mathcal O} \\
  r\coloneq R.x\bmod n,~\call{Check}{r\ne0} \\
  \Phi_i\coloneq\phi_i+\sum_{j\ne i}\psi_{j,i} \\
  V_i^a\coloneq\sum_{j\ne i}(y^a_{i,j}+z^a_{j,i}),\quad a\in\{\mathrm r,\mathrm x\} \\
  s_{1,i}\coloneq r_i\Phi_i+V_i^{\mathrm r} \\
  s^\circ_{0,i}\coloneq r(x'_i\Phi_i+V_i^{\mathrm x}) \\
  \mathrm{PreSignature}_i\coloneq(C,R,r,s_{1,i},s^\circ_{0,i},\phi_i) \\
  \pcreturn
\end{protocol}

RvoleRecv 失败时按 CheckP 处理发送方。
预签名连同上下文字段和一次性状态标记保存，其余临时矩阵与随机量销毁。
此时 $R$ 只是本方视图，完整点的一致性由第三轮检查。

\subsubsection{多方 Sign：第三轮}

采用 \href{https://eprint.iacr.org/2026/929.pdf}{2026/929，\S3}
的完整 nonce 点与签名符号检查。
读取预签名时先核对保存的消息、公钥、参与集合、参数与调用上下文；
发送签名份额之前原子地标记为已消费，成功、失败和崩溃恢复均不得再次使用。

\begin{protocol}[
  name=Sign,
  params={\mathrm{PreSignature}_i,m,Y',\mathbb S\mid C},
  returns={r,s},
  layout=multiline
]
  (C_i,R,r,s_{1,i},s^\circ_{0,i},\phi_i)\coloneq\mathrm{PreSignature}_i \\
  \call{Check}{C_i=C\ \land\ \mathrm{unused}(\mathrm{PreSignature}_i)} \\
  \call{Check}{R\ne\mathcal O\ \land\ r=R.x\bmod n\ \land\ r\ne0} \\
  s_{0,i}\coloneq s^\circ_{0,i}+m\phi_i \\
  \mathrm{Consume}(\mathrm{PreSignature}_i) \\
  \pcsend{C,R,s_{0,i},s_{1,i}}{\mathbb S\setminus\{i\}} \\
  \call{Check}{\forall j\ne i,~C_j=C\ \land\ R^{(j)}=R} \\
  D\coloneq\sum_{j\in\mathbb S}s_{1,j},~\call{Check}{D\ne0} \\
  s\coloneq D^{-1}\sum_{j\in\mathbb S}s_{0,j},~\call{Check}{s\ne0} \\
  \call{Check}{s^{-1}mG+s^{-1}rY'=R} \\
  \pcif\mathrm{signMode}=\texttt{low-S};\pcthen \\
  s\coloneq\min(s,n-s) \\
  \pcfi \\
  \call{Check}{\mathrm{EcdsaVerify}(m,Y',r,s)} \\
  \mathrm{Erase}(\mathrm{PreSignature}_i) \\
  \pcreturn
\end{protocol}

完整点比较包含纵坐标符号，符号等式在 low-S 规范化之前检查。
EcdsaVerify 验证 $1\le r,s<n$、公钥合法且非零，
并检查 $U=s^{-1}(mG+rY')\ne\mathcal O$ 及 $U.x\bmod n=r$。
所有失败出口也执行 Erase；Consume 的持久化先于网络发送。

\paragraph{现场零分享} 如选用 \href{https://eprint.iacr.org/2026/929.pdf}{2026/929，\S3}
的放宽零分享接口，令 $\mathrm{zeroMode}=\texttt{fresh}$，
Keygen 省去零分享种子的承诺、揭示和保存，其余初始化步骤照常执行。
PreSign 第一轮对每个 $i<j$，由 $i$ 采样 $a_{i,j}\gets\uniform{\mathbb Z_n}$ 并保密发送给 $j$。
第二轮开始时计算
\[
  \zeta_i=\sum_{j<i}a_{j,i}-\sum_{j>i}a_{i,j},\qquad
  x'_i=\lambda(i,\mathbb S)x_i+\zeta_i+\nabla x\,|\mathbb S|^{-1}.
\]
该模式将 $x'_i,Y_i$ 的计算移到第二轮生成 RVOLE 消息之前，轮数仍为三轮。
各方必须在 $C$ 中约定同一种模式，标量消息也绑定 $C$、双方编号与用途。

\subsubsection{两方独立预处理与两轮在线签名}

采用 \href{https://eprint.iacr.org/2026/976.pdf}{2026/976，Protocol 6.2}。
本节双方编号改为 $i\in\{0,1\}$，$j=1-i$，在线时提供有效且一致的加法私钥分片
$x'_0,x'_1$ 及公开分片 $Y_i=x'_iG$，共同公钥 $Y'=Y_0+Y_1$。
论文还要求兼容 Keygen 对恶意份额的提取；前述 Keygen 与该证明的对应仍需核实，
本节的代数编排不据此宣称完成了这项组合证明。

每份预处理使用唯一 $\mathtt{ppid}$，两个方向均独立初始化 Base OT 与 PPRF。
对每个方向，Receiver 调用 EndemicRecvInit，Sender 调用 EndemicSendInit，
双方同轮交换 $R^0,R^1$ 与 $V^0,V^1$，再分别调用 EndemicRecvFinish 和 EndemicSendFinish。
每 $k$ 个 OT 交给 PprfBuild，Sender 在下一轮发送各树的 $P$，Receiver 调用 PprfEval 校验。
会话及树编号、叶子保存和 $T^{\mathrm{init}}$ 的计算沿用 Keygen 的相应公式，
以 $\mathtt{ppid}$ 替换 $\mathtt{kgid}$；本次初始化不生成密钥份额。
这两轮初始化计入离线阶段，随后按下列协议执行两轮扩展与 RVOLE。

本节取 Variant II，$L=896,N_1=2,N_2=1$，有向实例号 $c_{i,j}$ 绑定 $\mathtt{ppid}$ 及双方角色。
$T_{i,j}$ 仍绑定该方向初始化和本批 $E_{i,j}$。
Nonce 揭示随机数采用本文的 $\lambda=256$ 比特口径。

\begin{protocol}[
  name=PreprocessTwoParty,
  params={A_{*,i,j},\delta_{*,j,i},A^*_{*,j,i},T^{\mathrm{init}}\mid\mathtt{ppid}},
  returns={\mathrm{State}_i},
  layout=multiline
]
  r_i\coloneq\uniform{\mathbb Z_n},~R_i\coloneq r_iG,~e_i^R\coloneq\uniform{\bits{\lambda}} \\
  C_i^R\coloneq H(\mathtt{ppid},\texttt{"nonce"},i,R_i,e_i^R) \\
  \boldsymbol\beta_{j,i}\coloneq\uniform{\bits{L}} \\
  (\rho^{\boldsymbol\beta_{j,i}}_{j,i},E_{j,i})\coloneq
    \call{SoftSpokenRecv}{A_{*,i,j},\boldsymbol\beta_{j,i}\mid c_{j,i},L} \\
  \pcsend{C_i^R,E_{j,i}}{j}\pccomment{离线扩展第一轮} \\
  (\rho^0_{i,j},\rho^1_{i,j})\coloneq
    \call{SoftSpokenSend}{\delta_{*,j,i},A^*_{*,j,i},E_{i,j}\mid c_{i,j},L} \\
  (w_{i,j},\alpha^0_{i,j},Q_{i,j})\coloneq \\
  \quad\call{RvoleSend}{\rho^0_{i,j},\rho^1_{i,j}\mid c_{i,j},T_{i,j},L,2,1} \\
  (\tilde r_i,\tilde x_i)\coloneq w_{i,j} \\
  g_{i,j}\coloneq\call{Gadget}{\mathrm{II},Q_{i,j}\mid c_{i,j},T_{i,j},L} \\
  a_i^b\coloneq-\sum_{\ell=0}^{L-1}g_{i,j,\ell}\alpha^0_{i,j,\ell,b},\quad b\in\{\mathrm r,\mathrm x\} \\
  \pcsend{Q_{i,j}}{j}\pccomment{离线扩展第二轮} \\
  t_{j,i}\coloneq\call{RvoleRecv}{\boldsymbol\beta_{j,i},\rho^{\boldsymbol\beta_{j,i}}_{j,i},Q_{j,i}\mid c_{j,i},T_{j,i},L,2,1} \\
  g_{j,i}\coloneq\call{Gadget}{\mathrm{II},Q_{j,i}\mid c_{j,i},T_{j,i},L} \\
  \beta_i\coloneq\langle g_{j,i},\boldsymbol\beta_{j,i}\rangle \\
  b_i^a\coloneq\sum_{\ell=0}^{L-1}g_{j,i,\ell}t_{j,i,\ell,a},\quad a\in\{\mathrm r,\mathrm x\} \\
  \mathrm{State}_i\coloneq(\beta_i,r_i,\tilde r_i,\tilde x_i,a_i^{\mathrm r},a_i^{\mathrm x},b_i^{\mathrm r},b_i^{\mathrm x}) \\
  \pcreturn
\end{protocol}

这里 $a,b$ 的下标表示份额持有人，因此
$a_i^{\mathrm r}+b_j^{\mathrm r}=\tilde r_i\beta_j$，
$a_i^{\mathrm x}+b_j^{\mathrm x}=\tilde x_i\beta_j$。
State 还保存对端承诺 $C_j^R$、本方揭示随机数 $e_i^R$、预处理编号、身份参数及待用标记。
成功压缩状态后销毁本批初始化秘密和临时矩阵；预处理失败也销毁整批材料。
此阶段不揭示 $R_i$，也不绑定实际签名密钥或消息。

在线时共同确定 $m,Y',Y_0,Y_1$，并检查 $Y'=Y_0+Y_1$ 及本方 $Y_i=x'_iG$。
在线上下文 $C_{\mathrm{on}}$ 绑定预处理编号、唯一签名号、双方身份、密钥标识、上述公开输入、参数及输出规则。
发送第一条在线消息之前，原子地将待用状态绑定 $C_{\mathrm{on}}$ 并标记已消费；
以后不能更换消息或密钥，取消、失败及崩溃恢复也不能将状态放回待用池。

\begin{protocol}[
  name=SignTwoParty,
  params={\mathrm{State}_i,x'_i,m,Y',Y_0,Y_1\mid C_{\mathrm{on}}},
  returns={r,s},
  layout=multiline
]
  \call{Check}{\mathrm{unused}(\mathrm{State}_i)\ \land\ Y'=Y_0+Y_1\ \land\ Y_i=x'_iG} \\
  \mathrm{Consume}(\mathrm{State}_i,C_{\mathrm{on}}) \\
  \gamma_i^{\mathrm r}\coloneq r_i-\tilde r_i,~\gamma_i^{\mathrm x}\coloneq x'_i-\tilde x_i \\
  \Gamma_i^{\mathrm r}\coloneq a_i^{\mathrm r}G,~\Gamma_i^{\mathrm x}\coloneq a_i^{\mathrm x}G \\
  R_i\coloneq r_iG \\
  \pcsend{C_{\mathrm{on}},\gamma_i^{\mathrm r},\gamma_i^{\mathrm x},\Gamma_i^{\mathrm r},\Gamma_i^{\mathrm x},R_i,e_i^R}{j}
    \pccomment{在线第一轮} \\
  \call{Check}{C_{\mathrm{on}}^{(j)}=C_{\mathrm{on}}} \\
  \call{Check}{H(\mathtt{ppid},\texttt{"nonce"},j,R_j,e_j^R)=C_j^R} \\
  z_i^{\mathrm r}\coloneq b_i^{\mathrm r}+\beta_i\gamma_j^{\mathrm r},~
    z_i^{\mathrm x}\coloneq b_i^{\mathrm x}+\beta_i\gamma_j^{\mathrm x} \\
  \call{Check}{z_i^{\mathrm r}G=\beta_iR_j-\Gamma_j^{\mathrm r}} \\
  \call{Check}{z_i^{\mathrm x}G=\beta_iY_j-\Gamma_j^{\mathrm x}} \\
  R\coloneq R_0+R_1,~\call{Check}{R\ne\mathcal O} \\
  r\coloneq R.x\bmod n,~\call{Check}{r\ne0} \\
  u_i\coloneq r_i\beta_i+a_i^{\mathrm r}+z_i^{\mathrm r} \\
  v_i\coloneq x'_i\beta_i+a_i^{\mathrm x}+z_i^{\mathrm x} \\
  w_i\coloneq m\beta_i+rv_i \\
  \pcsend{C_{\mathrm{on}},u_i,w_i}{j}\pccomment{在线第二轮} \\
  \call{Check}{C_{\mathrm{on}}^{(j)}=C_{\mathrm{on}}} \\
  D\coloneq u_0+u_1,~\call{Check}{D\ne0} \\
  s\coloneq(w_0+w_1)D^{-1},~\call{Check}{s\ne0} \\
  \call{Check}{\mathrm{EcdsaVerify}(m,Y',r,s)} \\
  \pcif\mathrm{signMode}=\texttt{low-S};\pcthen \\
  s\coloneq\min(s,n-s) \\
  \pcfi \\
  \mathrm{Erase}(\mathrm{State}_i) \\
  \pcreturn
\end{protocol}

所有失败出口均销毁该批一次性状态。
两方流程直接以 $\beta_i$ 作为盲化分片，满足
\[
  \begin{aligned}
    u_0+u_1&=(r_0+r_1)(\beta_0+\beta_1),\\
    v_0+v_1&=(x'_0+x'_1)(\beta_0+\beta_1).
  \end{aligned}
\]
因此其两轮在线阶段使用自己的消息格式，按 Protocol 6.2 验签后输出。

\end{document}

\renewcommand{\doctitle}{DKLS23 实施方式}
\begin{document}
\begin{center}
  {\LARGE\heitiLatin\heiti\mdseries\doctitle\par}
\end{center}
\vspace{1em}
% title: DKLS23 实施方式

\section{前置工作}

\subsection{朴素 OT-MtA}

\begin{protocol}[
  name=EcdsaSign,
  params={m,x \mid n,G},
  returns={r,s}
]
  k \coloneq \uniform{\mathbb{Z}_n^*} \\
  R \coloneq kG \\
  r \coloneq R.x \bmod n \pccomment{点 $R$ 的横坐标} \\
  \label{ecdsa-s}\tag{ecdsa-s}
  \boldsymbol{s} \coloneq k^{-1}(m+rx) \bmod n \\
\end{protocol}

为 \eqref{ecdsa-s} 引入乘项 $\gamma \ne 0 \pmod n$，则有
\begin{equation*}
  s = (k\gamma)^{-1}(m\gamma + rx\gamma) \bmod n.
\end{equation*}

朴素 MtA 协议：

\begin{protocol}[
  name=NaiveMtA,
  params={\party{S}: x_a,~ \party{R}: x_b \mid n},
  returns={\party{S}: y_a,~ \party{R}: y_b},
  st={y_a+y_b \equiv x_a x_b \pmod n}
]
  \pcparty{S}
  \pcfor t \in [0, \ell); \pcdo \\
  r_t \coloneq \uniform{\mathbb{Z}_n^*} \\
  m_{t,0} \coloneq r_t,~ m_{t,1} \coloneq r_t + x_a 2^t \bmod n \\
  \pcendfor{t} \\

  \pcparty{R}\sum_{t=0}^{\ell-1}d_t 2^t \coloneq x_b \\

  \pcparty{S,R}\pcfor t \in [0, \ell); \pcdo \\
    \party{S}: \varnothing,~ \party{R}: m_{t,d_t} \coloneq
    \mathrm{NaiveOT}\left(\party{S}: (m_{t,0}, m_{t,1}),~ \party{R}: d_t\right) \\
  \pcendfor{t} \\

  \pcparty{S}y_a \coloneq -\sum_t m_{t,0} \bmod n \\
  \pcparty{R}y_b \coloneq \sum_t m_{t, d_t} \bmod n \\
\end{protocol}

朴素二选一 OT 协议：

\begin{protocol}[
  name=NaiveOT,
  params={\party{S}: (m_0, m_1),~ \party{R}: d\in\bit},
  returns={\party{S}: \emptyset,~ \party{R}: m_d}
]
  \pcparty{S}
  \alpha \coloneq \uniform{\mathbb{Z}_n^*},~ A \coloneq \alpha G \\
  \pcsend{A}{\party{R}} \\

  \pcparty{R}
  \beta \coloneq \uniform{\mathbb{Z}_n^*},~ B \coloneq \beta G + dA \\
  \pcsend{B}{\party{S}} \\

  \pcparty{S}
  K_0 \coloneq \hash{\alpha B},~ K_1 \coloneq \hash{\alpha (B-A)} \\
  C_0 \coloneq \enc{K_0, m_0},~ C_1 \coloneq \enc{K_1, m_1} \\
  \pcsend{C_0, C_1}{\party{R}} \\
  \pcparty{R}
  K_d \coloneq \hash{\beta A} \\
  m_d \coloneq \dec{K_d, C_d} \\
  \pcreturn
\end{protocol}
\bigskip

\subsection{IKNP03}

制备持久化种子，用于后续 OT 协议。

\begin{protocol}[
  name=BaseOTSetup,
  params={\kappa},
  returns={\party{S}: (s, k^s),~ \party{R}: (k^0, k^1)}
]
  \pcparty{S,R}\pcfor i \in [0,\kappa); \pcdo \\
  \pcparty{R}
  k_i^0 \coloneq \uniform{\bits{\kappa}},~
  k_i^1 \coloneq \uniform{\bits{\kappa}} \\
  \pcparty{S}s_i \coloneq \uniform{\bit} \\
  \pcparty{S,R}
  \party{R}: \varnothing,~ \party{S}: k_i^{s_i} \coloneq
  \call{NaiveOT}{\party{R}: (k_i^0,k_i^1),~ \party{S}: s_i}
  \pccomment{注意角色互换}
  \\
  \pcendfor{i} \\
  \pcparty{S}
  s \coloneq \bigl(s_i \mid i\in[0,\kappa)\bigr),~
  k^s \coloneq \bigl(k_i^{s_i} \mid i\in[0,\kappa)\bigr) \\
  \pcparty{R}
  k^0 \coloneq \bigl(k_i^0 \mid i\in[0,\kappa)\bigr),~
  k^1 \coloneq \bigl(k_i^1 \mid i\in[0,\kappa)\bigr) \\
  \pcreturn
\end{protocol}
\bigskip

子协议 IknpExpand 意在用 $\kappa$ 个（对）种子，生成 $m$ 个（对）密钥原像。
Sender 获得成对的密钥原像 $q_{*,j}$ 与 $q_{*,j}\oplus s$；Receiver 获得所选密钥原像 $t^0$。
\pdfcomment[
  author={wojtek},
  subject={Test Comment}
]{
  Whatever.
}

\begin{protocol}[
  name=IknpExpand,
  params={\party{S}: (s,k^s),~ \party{R}: (k^0,k^1,b) \mid \mathtt{sid}, \mathtt{tag}, L},
  returns={\party{S}: (q,u),~ \party{R}: (t^0,u)},
  st={\forall j\in[0,L),~ q_{*,j}=t^0_{*,j}\oplus(b_j\odot s)},
  % layout=multiline
]
  \pcparty{R}\pcfor i \in [0,\kappa); \pcdo \\
  t^0_{i,*} \coloneq \prf{\mathtt{sid}, \mathtt{tag},k_i^0} \in \bits{L} \\
  t^1_{i,*} \coloneq \prf{\mathtt{sid}, \mathtt{tag},k_i^1} \in \bits{L} \\
  u_{i,*} \coloneq t^0_{i,*} \oplus t^1_{i,*} \oplus b \\
  \pcendfor{i} \\
  \pcsend{u}{\party{S}} \\
  \pcparty{S}\pcfor i \in [0,\kappa); \pcdo \\
  t_{i,*} \coloneq \prf{\mathtt{sid}, \mathtt{tag},k_i^{s_i}} \in \bits{L} \\
  q_{i,*} \coloneq t_{i,*} \oplus (s_i \odot u_{i,*}) \\
  \pcendfor{i} \\
  \pcreturn
  \pcnote{
    操作符 $\oplus$ 是按位异或，$\odot$ 是按位与；
    矩阵 $q,u,t^0$ 的形状均为 $\kappa\times L$。}
\end{protocol}

双方算出密钥，传输二选一消息。

\begin{protocol}[
  name=PayloadOT,
  params={\party{S}: (q,s,M^0,M^1),~ \party{R}: (t^0,b) \mid \mathtt{sid}, \mathtt{tag},m},
  returns={\party{S}: \varnothing,~ \party{R}: M^b},
  layout=multiline
]
  \pcparty{S}\pcfor j \in [0,m); \pcdo \\
  K_j^0 \coloneq \hash{\mathtt{sid}, \mathtt{tag},j,q_{*,j}} \\
  K_j^1 \coloneq \hash{\mathtt{sid}, \mathtt{tag},j,q_{*,j}\oplus s} \\
  C_j^0 \coloneq \enc{K_j^0,M_j^0},~
  C_j^1 \coloneq \enc{K_j^1,M_j^1} \\
  \pcsend{j,C_j^0,C_j^1}{\party{R}} \\
  \pcendfor{j} \\
  \pcparty{R}\pcfor j \in [0,m); \pcdo \\
  K_j^{b_j} \coloneq \hash{\mathtt{sid}, \mathtt{tag},j,t^0_{*,j}} \\
  M_j^{b_j} \coloneq \dec{K_j^{b_j},C_j^{b_j}} \\
  \pcendfor{j} \\
  \pcreturn
  \pcnote{
    $M^0\coloneq\bigl(M_j^0\mid j\in[0,m)\bigr)$、
    $M^1\coloneq\bigl(M_j^1\mid j\in[0,m)\bigr)$ 为两组输入消息，
    $M^b\coloneq\bigl(M_j^{b_j}\mid j\in[0,m)\bigr)$ 为所选消息族。}
\end{protocol}

把 BaseOTSetup、IknpExpand、PayloadOT 编排成为完整的 IKNP03 协议。

\begin{protocol}[
  name=IknpOT,
  params={\party{S}: (M^0,M^1),~ \party{R}: b \mid \mathtt{sid}, \mathtt{tag},m,\kappa},
  returns={\party{S}: \varnothing,~ \party{R}: M^b},
  layout=multiline
]
  \pcparty{S,R}
  \party{S}: (s,k^s),~ \party{R}: (k^0,k^1) \coloneq
  \call{BaseOTSetup}{\kappa} \\
  \party{S}: (q,u),~ \party{R}: (t^0,u) \coloneq \\
  \quad \mathrm{IknpExpand}
  \left(\party{S}: (s,k^s),~ \party{R}: (k^0,k^1,b) \mid \mathtt{sid}, \mathtt{tag},m\right) \\
  \party{S}: \varnothing,~ \party{R}: M^b \coloneq \\
  \quad \mathrm{PayloadOT}
  \left(\party{S}: (q,s,M^0,M^1),~ \party{R}: (t^0,b) \mid \mathtt{sid}, \mathtt{tag},m\right) \\
  \pcreturn
\end{protocol}

\subsection{KOS15}

KOS 在发送 Payload 密文之前，检查矩阵各行是否使用一致的选择向量。
$\sum$ 表示域 $\mathbb{F}$ 元素的加法，即按位异或。
点乘 “$\cdot$” 表示域 $\mathbb{F}=\mathrm{GF}(2^\kappa)$ 元素的乘法，
现在（2026年9月）的 AI 已经能描述这种乘法如何实施，这里不再赘述。
单比特 $b_j$ 视为域中的常数 $0$ 或 $1$，故 $b_j\cdot\chi_j$ 表示选取或丢弃 $\chi_j$。

\begin{protocol}[
  name=KosCheck,
  params={\party{S}: (q,s),~ \party{R}: (t^0,b) \mid \chi,L},
  returns={\party{S, R}: \varnothing},
  layout=multiline
]
  \pcparty{R}
  \tau \coloneq \sum_{j=0}^{L-1}t^0_{*,j}\cdot\chi_j \in \mathbb F \\
  \rho \coloneq \sum_{j=0}^{L-1}b_j\cdot\chi_j \in \mathbb F \\
  \pcsend{\tau,\rho}{\party{S}} \\
  \pcparty{S}
  Q \coloneq \sum_{j=0}^{L-1}q_{*,j}\cdot\chi_j \in \mathbb F \\
  \pcif Q \ne \tau\oplus(\rho\cdot s); \pcthen
  \pccomment{请读者自行验算等式两边，建议交给 AI。}\\
  \pcabort \\
  \pcfi \\
  \pcreturn
\end{protocol}

\begin{protocol}[
  name=KosOT,
  params={\party{S}: (s,k^s,M^0,M^1),~ \party{R}: (k^0,k^1,b) \mid \mathtt{sid}, \mathtt{tag},m,m_2},
  returns={\party{S}: \varnothing,~ \party{R}: M^b},
  layout=multiline
]
  \pcparty{S,R}m'\coloneq m+\kappa+m_2 \\
  \pcparty{R}
  \tilde b\coloneq\uniform{\bits{\kappa+m_2}},~
  b'\coloneq b\mathbin{\|}\tilde b \\
  \pcparty{S,R}
  \party{S}: (q,u),~ \party{R}: (t^0,u) \coloneq \\
  \quad \call{IknpExpand}
  {\party{S}: (s,k^s),~ \party{R}: (k^0,k^1,b') \mid \mathtt{sid}, \mathtt{tag},m'} \\
  \chi\coloneq\prf{
    \texttt{"kos-check"},\mathtt{sid}, \mathtt{tag},\kappa,m,m_2,u
  }\in\mathbb F^{m'} \\
  \call{KosCheck}{\party{S}: (q,s),~ \party{R}: (t^0,b') \mid \chi,m'} \\
  \party{S}: \varnothing,~ \party{R}: M^b \coloneq \\
  \quad \call{PayloadOT}{
    \party{S}: (q,s,M^0,M^1),~ \party{R}: (t^0,b)
    \mid \mathtt{sid}, \mathtt{tag},m
  } \\
  \pcreturn
\end{protocol}

\paragraph{陪跑比特} KosOT 给实际选择向量 $b$ 追加 $\kappa+m_2$ 个随机陪跑比特，
目的是隐藏检查响应 $\rho$ 中关于 $b$ 的信息，$m_2$ 是统计安全参数。
检查覆盖全部 $m'=m+\kappa+m_2$ 列，PayloadOT 只使用前 $m$ 列；
陪跑列没有对应消息，也无需派生消息密钥。
这一安排及其安全分析直接见
\href{https://www.iacr.org/archive/crypto2015/92160327/92160327.pdf\#page=11}
{Keller、Orsini、Scholl，CRYPTO 2015，\S3.3 的 Lemma 2 及恶意 Sender 情形的证明}。
论文中的统计安全参数 $s$ 对应本文的 $m_2$。

\section{DKLS23 本身}

以下固定素数阶群 $\langle G\rangle$，群阶为 $n$，标量运算均在 $\mathbb Z_n$ 中进行。
$\kappa=\lceil\log_2n\rceil$，$\lambda$ 是种子长度，$\lambda_c,\lambda_s$ 分别是计算与统计安全参数。
默认参数为
\[
  \kappa=\lambda=256,\qquad \lambda_c=\lambda_s=128,
  \qquad k=4,\quad q=2^k=16.
\]

$H$ 表示带用途标签的随机预言机接口，输入先作含类型、长度和顺序的无歧义编码。
未注明时输出为 $\bits{2\lambda_c}$；注明域、矩阵形状或比特长度时，按该类型扩展输出，
并把目标类型纳入域分离。到 $\mathbb Z_n$ 的采样使用 XOF 产生候选标量并拒绝大于等于 $n$ 的值。
种子扩展记作 $\prf{d,K}\in\bits{\ell}$，其中 $K$ 为秘密种子，$d$ 为公开域分离输入。
所有 $\mathtt{sid}$ 均绑定协议版本、参数、双方身份、角色与唯一实例号。
并列随机赋值表示各变量独立采样；数组按索引递增排列，无下标的数组名表示相应的完整有序族。
接收消息时先验证上下文、长度、标量和群元素的规范编码及群成员资格。
$\call{Check}{P}$ 表示检查命题 $P$，不成立则立即中止；初始化与签名的失败处理范围在编排节规定。

\subsection{Endemic OT}

采用 \href{https://eprint.iacr.org/2019/706.pdf}{Masny--Rindal，Fig.\ 8}
所用的 Endemic OT 消息与密钥公式。
将本地计算拆开，便于在 Keygen 中与其他消息合并发送。
单独运行时，Receiver 的 $R^0,R^1$ 与 Sender 的 $V^0,V^1$ 可以同轮发送。

\begin{protocol}[
  name=EndemicRecvInit,
  params={s\in\bits{\kappa}\mid\mathtt{sid}},
  returns={u,R^0,R^1}
]
  \pcfor a\in[0,\kappa);\pcdo \\
  u_a\coloneq\uniform{\mathbb Z_n},~w_a\coloneq\uniform{\bits{\lambda}} \\
  R_a^{1-s_a}\coloneq\mathrm{HashToGroup}(\mathtt{sid},\texttt{"ot-r"},a,w_a) \\
  R_a^{s_a}\coloneq u_aG-\mathrm{HashToGroup}(\mathtt{sid},\texttt{"ot-u"},s_a,a,R_a^{1-s_a}) \\
  \pcendfor{a} \\
  \pcreturn
\end{protocol}

\begin{protocol}[
  name=EndemicSendInit,
  params={\kappa},
  returns={v^0,v^1,V^0,V^1}
]
  \pcfor a\in[0,\kappa),~d\in\bit;\pcdo \\
  v_a^d\coloneq\uniform{\mathbb Z_n},~V_a^d\coloneq v_a^dG \\
  \pcendfor{a,d} \\
  \pcreturn
\end{protocol}

\begin{protocol}[
  name=EndemicSendFinish,
  params={v^0,v^1,R^0,R^1\mid\mathtt{sid}},
  returns={\rho^0,\rho^1}
]
  \pcfor a\in[0,\kappa),~d\in\bit;\pcdo \\
  U_a^d\coloneq R_a^d+\mathrm{HashToGroup}(\mathtt{sid},\texttt{"ot-u"},d,a,R_a^{1-d}) \\
  \rho_a^d\coloneq H(\mathtt{sid},\texttt{"ot-key"},d,a,v_a^dU_a^d)\in\bits{\lambda} \\
  \pcendfor{a,d} \\
  \pcreturn
\end{protocol}

\begin{protocol}[
  name=EndemicRecvFinish,
  params={s,u,V^0,V^1\mid\mathtt{sid}},
  returns={\rho^s}
]
  \pcfor a\in[0,\kappa);\pcdo \\
  \rho_a^{s_a}\coloneq H(\mathtt{sid},\texttt{"ot-key"},s_a,a,u_aV_a^{s_a})\in\bits{\lambda} \\
  \pcendfor{a} \\
  \pcreturn
\end{protocol}

这里 $\mathrm{HashToGroup}$ 使用离散对数未知的点哈希，不能实现为 $H(\cdots)G$；
$w_a$ 新鲜且保密。正确性来自 $U_a^{s_a}=u_aG$。
Endemic OT 允许恶意方影响自己所得密钥的分布，后续组合须保留这种接口条件，
不能把它直接当作双方输出均不可操纵的随机 OT。

\subsection{PPRF}

采用 \href{https://eprint.iacr.org/2022/192.pdf}{Roy，SoftSpokenOT，Fig.\ 13--14}
的建树、打孔与叶子校验步骤。每棵树消耗 $k$ 个 Base OT；
完整叶子族记为 $A$，缺少位置 $\delta$ 的叶子族记为 $A^*$。
$\mathtt{tid}$ 绑定初始化会话及树编号；本节 PRF 的域分离输入均包含 $\mathtt{tid}$。

\begin{protocol}[
  name=PprfBuild,
  params={\rho^0,\rho^1\mid\mathtt{tid},k},
  returns={A,P},
  layout=multiline
]
  T^1_0\coloneq\rho_0^0,~T^1_1\coloneq\rho_0^1 \\
  \pcfor a\in[1,k);\pcdo \\
  \pcfor z\in[0,2^a);\pcdo \\
  T^{a+1}_{2z}\mathbin{\|}T^{a+1}_{2z+1}
    \coloneq\prf{\mathtt{tid},\texttt{"tree"},T^a_z}\in\bits{2\lambda} \\
  \pcendfor{z} \\
  t_a^d\coloneq\rho_a^d\oplus\bigoplus_{z=0}^{2^a-1}T^{a+1}_{2z+d},\quad d\in\bit \\
  \pcendfor{a} \\
  A\coloneq(T^k_z\mid z\in[0,q)) \\
  \tilde s_z\coloneq\prf{\mathtt{tid},\texttt{"leaf"},T^k_z}\in\bits{2\lambda},\quad z\in[0,q) \\
  \tilde t\coloneq\bigoplus_{z=0}^{q-1}\tilde s_z \\
  \tilde s\coloneq H(\mathtt{tid},\texttt{"leaves"},(\tilde s_z\mid z\in[0,q)))\in\bits{2\lambda} \\
  P\coloneq(t^0,t^1,\tilde t,\tilde s) \\
  \pcreturn
\end{protocol}

\begin{protocol}[
  name=PprfEval,
  params={s,\rho^s,P\mid\mathtt{tid},k},
  returns={\delta,A^*}
]
  (t^0,t^1,\tilde t,\tilde s)\coloneq P \\
  y_1\coloneq1-s_0,~T^1_{s_0}\coloneq\rho_0^{s_0} \\
  \pcfor a\in[1,k);\pcdo \\
  \pcfor z\in[0,2^a)\setminus\{y_a\};\pcdo \\
  T^{a+1}_{2z}\mathbin{\|}T^{a+1}_{2z+1}
    \coloneq\prf{\mathtt{tid},\texttt{"tree"},T^a_z}\in\bits{2\lambda} \\
  \pcendfor{z} \\
  T^{a+1}_{2y_a+s_a}\coloneq t_a^{s_a}\oplus\rho_a^{s_a}
    \oplus\bigoplus_{z\ne y_a}T^{a+1}_{2z+s_a} \\
  y_{a+1}\coloneq2y_a+1-s_a \\
  \pcendfor{a} \\
  \delta\coloneq y_k,~A^*\coloneq(T^k_z\mid z\in[0,q)\setminus\{\delta\}) \\
  \pcfor z\in[0,q)\setminus\{\delta\};\pcdo \\
  \tilde s_z\coloneq\prf{\mathtt{tid},\texttt{"leaf"},T^k_z}\in\bits{2\lambda} \\
  \pcendfor{z} \\
  \tilde s_\delta\coloneq\tilde t\oplus\bigoplus_{z\ne\delta}\tilde s_z \\
  \call{Check}{H(\mathtt{tid},\texttt{"leaves"},(\tilde s_z\mid z\in[0,q)))=\tilde s} \\
  \pcreturn
\end{protocol}

Sender 发送 $P$，Receiver 调用 PprfEval；校验通过后才保存初始化结果。
Base OT 选择位与打孔位置的关系为
\[
  \delta=\sum_{a=0}^{k-1}(1-s_a)2^{k-1-a}.
\]
后文统一按高位在前编码 $\delta$，即 $\operatorname{bit}_a(\delta)=1-s_a$。
每个有向初始化还保存 Base OT 与 PPRF 公共消息的规范 transcript 摘要 $T_{\mathrm{init}}$，
包含会话、身份、角色、参数及版本；消息按协议规定的顺序编码，不按本地到达顺序编码。

\subsection{SoftSpokenOT}

完整树持有者成为扩展 OT 的 Receiver，打孔树持有者成为扩展 OT 的 Sender。
每批使用 $h\in[0,\kappa/k)$ 的树，扩展出 $L$ 个随机 OT。
以下将 \href{https://eprint.iacr.org/2022/192.pdf}{SoftSpokenOT}
的展开与一致性检查合为一条 Receiver 消息，使用 05 中的随机预言机挑战实现。

令 $\mathbb F_c=\mathrm{GF}(2^\sigma)$，固定其多项式基，默认 $\sigma=128$。
$L$ 是 $\sigma$ 的倍数，$B=L/\sigma$，$L'=L+\sigma$。
将 $x\in\bits{L'}$ 分为 $B+1$ 个域元素 $x^{[0]},\ldots,x^{[B]}$，定义
\[
  h_\chi(x)\coloneq x^{[B]}+\sum_{b=0}^{B-1}\chi_bx^{[b]}
  \in\mathbb F_c.
\]
本节域加法为异或；$\odot$ 表示由单比特控制的取值或置零。

\begin{protocol}[
  name=SoftSpokenRecv,
  params={A,\boldsymbol\beta\in\bits{L}\mid\mathtt{sid},L},
  returns={\rho^{\boldsymbol\beta},E},
  layout=multiline
]
  \beta^{\mathrm{ext}}\coloneq\uniform{\bits{\sigma}},~
    \hat\beta\coloneq\boldsymbol\beta\mathbin{\|}\beta^{\mathrm{ext}} \\
  \pcfor h\in[0,\kappa/k);\pcdo \\
  r_{h,x}\coloneq\prf{\mathtt{sid},h,\texttt{"expand"},A_{h,x}}\in\bits{L'},\quad x\in[0,q) \\
  u_h\coloneq\hat\beta\oplus\bigoplus_{x=0}^{q-1}r_{h,x} \\
  \pcfor a\in[0,k);\pcdo \\
  v_{hk+a,*}\coloneq\bigoplus_{x=0}^{q-1}\operatorname{bit}_a(x)\odot r_{h,x} \\
  \pcendfor{a} \\
  \pcendfor{h} \\
  \chi\coloneq H(\mathtt{sid},\texttt{"ss-check"},u)\in\mathbb F_c^B \\
  \tau_a\coloneq h_\chi(v_{a,*}),\quad a\in[0,\kappa) \\
  \tilde\beta\coloneq h_\chi(\hat\beta),~E\coloneq(u,\tau,\tilde\beta) \\
  \rho_j^{\beta_j}\coloneq H(\mathtt{sid},\texttt{"ss-key"},j,v_{*,j})\in\bits{\lambda},\quad j\in[0,L) \\
  \pcreturn
\end{protocol}

\begin{protocol}[
  name=SoftSpokenSend,
  params={\delta,A^*,E\mid\mathtt{sid},L},
  returns={\rho^0,\rho^1}
]
  (u,\tau,\tilde\beta)\coloneq E \\
  \chi\coloneq H(\mathtt{sid},\texttt{"ss-check"},u)\in\mathbb F_c^B \\
  \pcfor h\in[0,\kappa/k);\pcdo \\
  \pcfor x\in[0,q)\setminus\{\delta_h\};\pcdo \\
  r_{h,x}\coloneq\prf{\mathtt{sid},h,\texttt{"expand"},A^*_{h,x}}\in\bits{L'} \\
  \pcendfor{x} \\
  \pcfor a\in[0,k);\pcdo \\
  \Delta_{hk+a}\coloneq\operatorname{bit}_a(\delta_h) \\
  w_{hk+a,*}\coloneq\bigoplus_{x\ne\delta_h}\operatorname{bit}_a(\delta_h\oplus x)\odot r_{h,x} \\
  w_{hk+a,*}\coloneq w_{hk+a,*}\oplus(\Delta_{hk+a}\odot u_h) \\
  \call{Check}{h_\chi(w_{hk+a,*})=\tau_{hk+a}\oplus(\Delta_{hk+a}\odot\tilde\beta)} \\
  \pcendfor{a} \\
  \pcendfor{h} \\
  \pcfor j\in[0,L);\pcdo \\
  \rho_j^0\coloneq H(\mathtt{sid},\texttt{"ss-key"},j,w_{*,j})\in\bits{\lambda} \\
  \rho_j^1\coloneq H(\mathtt{sid},\texttt{"ss-key"},j,w_{*,j}\oplus\Delta)\in\bits{\lambda} \\
  \pcendfor{j} \\
  \pcreturn
\end{protocol}

Receiver 发送 $E$，Sender 验证后才派生输出密钥。
诚实执行满足 $w_{*,j}=v_{*,j}\oplus(\hat\beta_j\odot\Delta)$；
最后 $\sigma$ 列只参与检查，输出仅使用前 $L$ 列。
每批保存
\[
  T_{\mathrm{OT}}\coloneq H(\mathtt{sid},\texttt{"ot-transcript"},T_{\mathrm{init}},E).
\]
初始化摘要不能由当前扩展消息代替。
挑战的具体 transcript 绑定沿用
\href{https://eprint.iacr.org/2026/976.pdf}{2026/976，\S4.5} 的组合要求；
Endemic OT、PPRF、扩展 OT 的组合还须满足该处的输入绑定条件。

\subsection{RVOLE}

保留 \href{https://eprint.iacr.org/2026/976.pdf}{2026/976，\S4}
的完整一致性检查和随机输入接口，指定输入通过偏移接入。
功能负载数 $N_1=2$，默认校验负载数 $N_2=1$；
协议暂按一般的 $N_1,N_2$ 书写，令 $N=N_1+N_2$。
Receiver 为每个实例独立采样 $\boldsymbol\beta\in\bits{L}$，
经 SoftSpokenOT 获得 $\rho^{\boldsymbol\beta}$，Sender 获得 $\rho^0,\rho^1$。
二者使用同一个 $T_{\mathrm{OT}}$。

定义域上的密钥扩展
\[
  \alpha^d_{j,a}\coloneq\mathrm{PRF}_{\rho_j^d}
  \circ H(\mathtt{sid},\texttt{"ot-pad"},j,a)\in\mathbb Z_n.
\]
这里的 PRF 输出也按前述拒绝采样规则产生标量。
RvoleSend 先生成消息及逐行份额，RvoleRecv 验证后产生逐行份额；
最后的 gadget 聚合单独列出，以兼容两种生成时机。

\begin{protocol}[
  name=RvoleSend,
  params={\rho^0,\rho^1\mid\mathtt{sid},T_{\mathrm{OT}},L,N_1,N_2},
  returns={w,\alpha^0,Q},
  layout=multiline
]
  w_a\coloneq\uniform{\mathbb Z_n},\quad a\in[0,N) \\
  \pcfor j\in[0,L),~a\in[0,N);\pcdo \\
  \alpha^d_{j,a}\coloneq\mathrm{PRF}_{\rho_j^d}\circ
    H(\mathtt{sid},\texttt{"ot-pad"},j,a)\in\mathbb Z_n,\quad d\in\bit \\
  \tilde a_{j,a}\coloneq\alpha^0_{j,a}-\alpha^1_{j,a}+w_a \\
  \pcendfor{j,a} \\
  \Theta\coloneq H(\mathtt{sid},\texttt{"challenge"},T_{\mathrm{OT}},\tilde a)
    \in\mathbb Z_n^{N_2\times N_1} \\
  \pcfor a\in[N_1,N);\pcdo \\
  \eta_a\coloneq w_a+\sum_{b=0}^{N_1-1}\theta_{a,b}w_b \\
  \mu_{j,a}\coloneq\alpha^0_{j,a}+\sum_{b=0}^{N_1-1}\theta_{a,b}\alpha^0_{j,b},\quad j\in[0,L) \\
  \pcendfor{a} \\
  \mu\coloneq H(\mathtt{sid},\texttt{"check"},(\mu_{j,a}\mid j\in[0,L),~a\in[N_1,N))) \\
  Q\coloneq(\tilde a,\eta,\mu) \\
  w\coloneq(w_b\mid b\in[0,N_1)) \\
  \pcreturn
\end{protocol}

\begin{protocol}[
  name=RvoleRecv,
  params={\boldsymbol\beta,\rho^{\boldsymbol\beta},Q\mid\mathtt{sid},T_{\mathrm{OT}},L,N_1,N_2},
  returns={t},
  layout=multiline
]
  (\tilde a,\eta,\mu)\coloneq Q \\
  \pcfor j\in[0,L),~a\in[0,N);\pcdo \\
  \alpha^{\beta_j}_{j,a}\coloneq\mathrm{PRF}_{\rho_j^{\beta_j}}\circ
    H(\mathtt{sid},\texttt{"ot-pad"},j,a)\in\mathbb Z_n \\
  t_{j,a}\coloneq\alpha^{\beta_j}_{j,a}+\beta_j\tilde a_{j,a} \\
  \pcendfor{j,a} \\
  \Theta\coloneq H(\mathtt{sid},\texttt{"challenge"},T_{\mathrm{OT}},\tilde a)
    \in\mathbb Z_n^{N_2\times N_1} \\
  \pcfor j\in[0,L),~a\in[N_1,N);\pcdo \\
  \hat\mu_{j,a}\coloneq t_{j,a}+\sum_{b=0}^{N_1-1}\theta_{a,b}t_{j,b}-\beta_j\eta_a \\
  \pcendfor{j,a} \\
  \hat\mu\coloneq H(\mathtt{sid},\texttt{"check"},(\hat\mu_{j,a}\mid j\in[0,L),~a\in[N_1,N))) \\
  \call{Check}{\hat\mu=\mu} \\
  \pcreturn
\end{protocol}

矩阵与族均按索引递增编码；$\Theta$ 的行按 $a\in[N_1,N)$ 编号。
检查通过后再使用功能负载，校验负载不输出给签名层。
诚实执行时 $t_{j,a}=\alpha^0_{j,a}+\beta_jw_a$，因此检查中所有 $\beta_j$ 项抵消。
这条正确性等式不替代恶意安全分析，后者还要求相应的随机 OT 或具体 transcript 输入绑定性质。

\subsection{gadget 合成与分解}

gadget 将较长的选择串合成为域上的标量，同时聚合双方的逐行乘法份额。
采用 \href{https://eprint.iacr.org/2026/976.pdf}{2026/976，\S4--5}
中适用于两种编排的方案：多方流程取 III，两方独立预处理取 II。
二者使用同一组 RVOLE 消息公式，只改变 $g$ 的产生时机与 $L$。

\begin{protocol}[
  name=Gadget,
  params={v,Q\mid\mathtt{sid},T_{\mathrm{OT}},L},
  returns={g\in\mathbb Z_n^L}
]
  \pcif v=\mathrm{III};\pcthen \\
  g\coloneq H(\mathtt{sid},\texttt{"gadget-III"})\in\mathbb Z_n^L \\
  \pcelse \\
  \call{Check}{v=\mathrm{II}\ \land\ Q\ne\bot} \\
  D_Q\coloneq H(\mathtt{sid},\texttt{"rvole-msg"},Q) \\
  g\coloneq H(\mathtt{sid},\texttt{"gadget-II"},T_{\mathrm{OT}},D_Q)\in\mathbb Z_n^L \\
  \pcfi \\
  \pcreturn
\end{protocol}

III 在实例开始时调用，允许输入 $Q=\bot$。
II 必须等 RvoleSend 的最终消息 $Q$ 固定后调用；Receiver 先验证 $Q$，再输出聚合结果。
III 的 $T_{\mathrm{OT}}$ 尚未产生时也可传入 $\bot$，该分支不使用它。

\begin{protocol}[
  name=RvoleAggregate,
  params={\party{S}: (w,\alpha^0),~\party{R}: (\boldsymbol\beta,t)\mid g,L,N_1},
  returns={\party{S}: (w,y),~\party{R}: (\beta,z)},
  st={\forall a\in[0,N_1),~y_a+z_a=w_a\beta},
  layout=multiline
]
  \pcparty{R}\beta\coloneq\sum_{j=0}^{L-1}g_j\beta_j \\
  \pcparty{S,R}\pcfor a\in[0,N_1);\pcdo \\
  \pcparty{S}y_a\coloneq-\sum_{j=0}^{L-1}g_j\alpha^0_{j,a} \\
  \pcparty{R}z_a\coloneq\sum_{j=0}^{L-1}g_jt_{j,a} \\
  \pcendfor{a} \\
  \pcreturn
\end{protocol}

指定 Sender 输入 $a$ 时，使用以下偏移；$w$ 必须是本次新鲜采样的随机量。
Sender 若已经知道 $a$，可将偏移与 $Q$ 同轮发送。

\begin{protocol}[
  name=RvoleShift,
  params={\party{S}: (a,w,y),~\party{R}: (\beta,z)},
  returns={\party{S}: y,~\party{R}: z},
  st={y+z=a\beta},
  layout=multiline
]
  \pcparty{S}\delta\coloneq a-w \\
  \pcsend{\delta}{\party{R}} \\
  \pcparty{R}z\coloneq z+\beta\delta \\
  \pcreturn
\end{protocol}

所谓分解，是在生成时保留 $\boldsymbol\beta$ 以及关系 $\beta=\langle g,\boldsymbol\beta\rangle$；
此处没有把任意给定标量高效逆解成选择串的步骤。
对于指定的 Sender 输入，用 RvoleShift 接入即可。

按上述文献的充分条件，II 要求 $L\ge\kappa+2\lambda_c+3\lambda_s$；III 要求
\[
  L\ge\kappa+7\lambda_s+2\log_2(\lambda_s+1)
     +2\lambda_s\log_2(eL/\lambda_s).
\]
取单个校验负载还要求 $\log_2n\ge\lambda_c+2\log_2L$。
对前述默认参数，按 $\sigma=128$ 对齐后使用
\[
  \begin{array}{c|rrrr}
    \text{用途}&L&L'&B&N_2\\\hline
    \text{多方，III}&2688&2816&21&1\\
    \text{两方独立预处理，II}&896&1024&7&1
  \end{array}
\]
两种方案都使用 $256$ 个 Base OT、$64$ 棵深度为 $4$ 的树。
这些是相应 RVOLE 变体的参数条件，整套签名的安全结论还取决于底层实现与组合前提。

\subsection{编排}

参与方集合 $\mathbb K\subset\mathbb Z_n^*$，各编号互异且非零；
阈值为 $t$，签名集合 $\mathbb S\subseteq\mathbb K$ 满足 $|\mathbb S|\ge t$。
标量与秘密种子通过认证保密信道传输，广播消息须有一致的发送者身份与会话标识。
下列伪代码在每方 $i$ 执行，$j$ 始终表示其他参与方；
每轮收齐所需消息并通过检查后，才进入下一轮。

树的下标统一按 PPRF Sender、Receiver 排列：$A_{h,i,j}$ 由 $i$ 持有，
对应的 $(\delta_{h,i,j},A^*_{h,i,j})$ 由 $j$ 持有。
扩展 OT 和 RVOLE 则按各自的 Sender、Receiver 排列下标。
因此，有向 RVOLE $i\to j$ 使用 $j$ 的完整树 $A_{h,j,i}$ 和 $i$ 的打孔树 $A^*_{h,j,i}$，
其 SoftSpoken 消息 $E_{i,j}$ 由 $j$ 发给 $i$，RVOLE 消息 $Q_{i,j}$ 由 $i$ 发给 $j$。

\paragraph{失败与状态} Keygen 中任何检查失败均中止初始化，销毁临时秘密，不输出 Keyshare。
签名中任何失败均废弃本次预签名，重试使用新会话、新随机数和新选择串。
RVOLE、逐对 $\Gamma$、SoftSpoken 或 nonce 承诺检查失败时，还须持久停用相应对端，
取消所有共享该 OT 状态且涉及该对端的并发会话与待用预签名。
这个停用记录须在共享状态的进程和设备间生效，更换会话号不能清除它。

下文用 $\call{CheckP}{P,j}$ 表示检查失败时停用对端 $j$，其余 Check 只中止当前执行。
对无法定位的承诺视图或聚合公钥不一致，沿用 07 的保守策略，停用 $\mathbb S\setminus\{i\}$；
这是停止交互的范围，不是公开归责结论。
收到失败通知或因超时取消时，只终止相应执行，不仅凭通知或超时停用某方。
失败通知绑定会话且不含秘密；成功、失败、取消后的临时材料都按相同生命周期销毁。
RVOLE 与 $\Gamma$ 失败后的停止交互要求见
\href{https://eprint.iacr.org/2023/765.pdf}{DKLS23，Protocol 3.6，第 8--10 步}；
共享 OT 状态的限制见 \href{https://eprint.iacr.org/2026/976.pdf}{2026/976，\S2.1}。

\subsubsection{Keygen}

这里保留 07 所用的 Feldman 承诺与 DLog 知识证明版本。
它的实现来源尚未逐项对应论文，不能套用
\href{https://eprint.iacr.org/2023/765.pdf}{DKLS23，Protocol 7.1}
对 relaxed Keygen 的证明。以下将其实际计算、校验及 OT 初始化写全。

对公开点 $X=xG$ 的知识证明采用 Schnorr 形式。
证明上下文 $c$ 绑定 Keygen 会话、证明者及系数编号或最终份额用途。
证明者取 $a\gets\uniform{\mathbb Z_n}$，令
\[
  A=aG,\qquad e=H(c,\texttt{"dlog"},X,A)\in\mathbb Z_n,
  \qquad \pi=(A,z),\quad z=a+ex.
\]
$\mathrm{DLogProve}(x\mid c)$ 返回 $\pi$；
$\mathrm{DLogVerify}(X,\pi\mid c)$ 在验证编码后，重算 $e$ 并返回 $zG=A+eX$ 的真假。
多项式点承诺的求值和 Lagrange 系数记为
\[
  F_i(z)\coloneq\sum_{a=0}^{t-1}z^aF_{i,a},\qquad
  \lambda(i,J)\coloneq\prod_{j\in J\setminus\{i\}}\frac{-j}{i-j}.
\]

\begin{protocol}[
  name=Keygen,
  params={i\in\mathbb K\mid\mathtt{kgid},\mathbb K,t},
  returns={\mathrm{Keyshare}_i},
  layout=multiline
]
  c_{i,a}\coloneq\uniform{\mathbb Z_n},~F_{i,a}\coloneq c_{i,a}G,\quad a\in[0,t) \\
  P_i(X)\coloneq\sum_{a=0}^{t-1}c_{i,a}X^a \\
  e_i\coloneq\uniform{\bits{\lambda}},~C_i\coloneq H(\mathtt{kgid},\texttt{"feldman"},i,F_i,e_i) \\
  \pcfor j\in\mathbb K\setminus\{i\};\pcdo \\
  \epsilon^\circ_{i,j},\nu_{i,j}\coloneq\uniform{\bits{\lambda}} \\
  C^\epsilon_{i,j}\coloneq H(\mathtt{kgid},\texttt{"zero-seed"},i,j,\epsilon^\circ_{i,j},\nu_{i,j}) \\
  \pcsend{C_i,C^\epsilon_{i,j}}{j}\pccomment{第一轮} \\
  \pcendfor{j} \\
  \pi_{i,a}\coloneq\mathrm{DLogProve}(c_{i,a}\mid\mathtt{kgid},i,\texttt{"coefficient"},a),\quad a\in[0,t) \\
  \pcfor j\in\mathbb K\setminus\{i\};\pcdo \\
  b_{j,i}\coloneq H(\mathtt{kgid},\texttt{"base"},j,i) \\
  s_{j,i}\coloneq\uniform{\bits{\kappa}} \\
  (u_{j,i},R^0_{j,i},R^1_{j,i})\coloneq\call{EndemicRecvInit}{s_{j,i}\mid b_{j,i}} \\
  \pcsend{F_i,e_i,\pi_i,R^0_{j,i},R^1_{j,i},\epsilon^\circ_{i,j},\nu_{i,j}}{j}
    \pccomment{第二轮} \\
  \pcendfor{j} \\
  \pcfor j\in\mathbb K\setminus\{i\};\pcdo \\
  \call{Check}{C_j=H(\mathtt{kgid},\texttt{"feldman"},j,F_j,e_j)} \\
  \call{Check}{C^\epsilon_{j,i}=H(\mathtt{kgid},\texttt{"zero-seed"},j,i,\epsilon^\circ_{j,i},\nu_{j,i})} \\
  \call{Check}{\forall a\in[0,t),~\mathrm{DLogVerify}(F_{j,a},\pi_{j,a}\mid\mathtt{kgid},j,\texttt{"coefficient"},a)} \\
  \epsilon_{i,j}\coloneq\epsilon^\circ_{i,j}\oplus\epsilon^\circ_{j,i} \\
  \pcendfor{j} \\
  F\coloneq\sum_{j\in\mathbb K}F_j,~d_{i,j}\coloneq P_i(j),\quad j\in\mathbb K \\
  \pcfor j\in\mathbb K\setminus\{i\};\pcdo \\
  b_{i,j}\coloneq H(\mathtt{kgid},\texttt{"base"},i,j) \\
  (v^0,v^1,V^0_{i,j},V^1_{i,j})\coloneq\call{EndemicSendInit}{\kappa} \\
  (\rho^0_{i,j},\rho^1_{i,j})\coloneq\call{EndemicSendFinish}{v^0,v^1,R^0_{i,j},R^1_{i,j}\mid b_{i,j}} \\
  \pcfor h\in[0,\kappa/k);\pcdo \\
  \mathtt{tid}_{h,i,j}\coloneq H(b_{i,j},\texttt{"pprf"},h) \\
  (A_{h,i,j},P_{h,i,j})\coloneq \\
  \quad\call{PprfBuild}{\rho^0_{i,j}[hk:(h+1)k],\rho^1_{i,j}[hk:(h+1)k]\mid\mathtt{tid}_{h,i,j},k} \\
  \pcendfor{h} \\
  \pcsend{d_{i,j},F,V^0_{i,j},V^1_{i,j},(P_{h,i,j}\mid h\in[0,\kappa/k))}{j}
    \pccomment{第三轮} \\
  \pcendfor{j} \\
  \pcfor j\in\mathbb K\setminus\{i\};\pcdo \\
  \call{Check}{F^{(j)}=F\ \land\ d_{j,i}G=F_j(i)} \\
  \rho^{s_{j,i}}_{j,i}\coloneq\call{EndemicRecvFinish}{s_{j,i},u_{j,i},V^0_{j,i},V^1_{j,i}\mid b_{j,i}} \\
  \pcfor h\in[0,\kappa/k);\pcdo \\
  \mathtt{tid}_{h,j,i}\coloneq H(b_{j,i},\texttt{"pprf"},h) \\
  (s',\rho')\coloneq(s_{j,i}[hk:(h+1)k],\rho^{s_{j,i}}_{j,i}[hk:(h+1)k]) \\
  (\delta_{h,j,i},A^*_{h,j,i})\coloneq\call{PprfEval}{s',\rho',P_{h,j,i}\mid\mathtt{tid}_{h,j,i},k} \\
  \pcendfor{h} \\
  T^{\mathrm{init}}_{i,j}\coloneq H(b_{i,j},\texttt{"init"},R^0_{i,j},R^1_{i,j},V^0_{i,j},V^1_{i,j},P_{*,i,j}) \\
  T^{\mathrm{init}}_{j,i}\coloneq H(b_{j,i},\texttt{"init"},R^0_{j,i},R^1_{j,i},V^0_{j,i},V^1_{j,i},P_{*,j,i}) \\
  \pcendfor{j} \\
  x_i\coloneq\sum_{j\in\mathbb K}d_{j,i},~Y\coloneq F(0),~S_i\coloneq x_iG \\
  \pi_i\coloneq\mathrm{DLogProve}(x_i\mid\mathtt{kgid},i,\texttt{"share"}) \\
  \pcsend{S_i,\pi_i}{\mathbb K\setminus\{i\}}\pccomment{第四轮} \\
  \call{Check}{\forall j\ne i,~\mathrm{DLogVerify}(S_j,\pi_j\mid\mathtt{kgid},j,\texttt{"share"})} \\
  \call{Check}{Y\ne\mathcal O\ \land\ Y=\sum_{j\in\mathbb K}\lambda(j,\mathbb K)S_j} \\
  \mathrm{Keyshare}_i\coloneq(i,x_i,Y,\mathbb K,t,\mathtt{kgid},A_{*,i,*},\delta_{*,*,i},A^*_{*,*,i},T^{\mathrm{init}},\epsilon_{i,*}) \\
  \pcreturn
\end{protocol}

切片 $v[a:b]$ 表示保留 $a\le j<b$ 的分量，并从零重新编号供子协议使用。
$F^{(j)}$ 是 $j$ 发来的聚合承诺，$P_{*,i,j}$ 按树编号排序。
Keyshare 另外保存协议参数、密钥标识和初始为空的持久停用记录；
输出后销毁多项式系数、份额传输副本、Base OT 临时指数及不再需要的中间种子。

\subsubsection{多方 PreSign：前两轮}

第一轮之前固定消息哈希 $m$、派生偏移 $\nabla x$ 及参与集合 $\mathbb S$。
PreSign 只表示三轮签名的前两轮；其输出已经绑定消息和公钥。
多方流程采用 Variant III，参数为 $L=2688,N_1=2,N_2=1$。
各方先检查参与集合、Keyshare 状态及停用记录，并固定
\[
  \begin{aligned}
    Y'&\coloneq Y+\nabla xG,\\
    C&\coloneq H(\texttt{"sign"},\mathtt{sid},\mathrm{keyid},\mathrm{sort}(\mathbb S),Y',m,\\
      &\hspace{4em}\mathrm{params},\mathrm{zeroMode},\mathrm{signMode}).
  \end{aligned}
\]
所有后续消息都携带 $C$，输入不符则中止。
有向子实例号不得重复，且必须绑定双方角色与调用编号。
$g_{i,j}$ 由 Gadget 的 III 分支在该实例开始时产生；收到扩展消息后计算 $T_{i,j}$。
相应公式为
\[
  \begin{aligned}
    c_{i,j}&\coloneq H(C,\texttt{"rvole"},i,j,\mathrm{callid}),\\
    T_{i,j}&\coloneq H(c_{i,j},\texttt{"ot-transcript"},T^{\mathrm{init}}_{j,i},E_{i,j}).
  \end{aligned}
\]

\begin{protocol}[
  name=PreSignRoundOne,
  params={\mathrm{Keyshare}_i,m,\nabla x,\mathbb S\mid C},
  returns={\mathrm{St}_i},
  layout=multiline
]
  r_i,\phi_i\coloneq\uniform{\mathbb Z_n^*},~R_i\coloneq r_iG \\
  \pcfor j\in\mathbb S\setminus\{i\};\pcdo \\
  I\coloneq\min(i,j),~J\coloneq\max(i,j) \\
  \zeta_{i,j}\coloneq H(C,\texttt{"zero-mask"},I,J,\epsilon_{i,j})\in\mathbb Z_n \\
  \pcendfor{j} \\
  \zeta_i\coloneq\sum_{j<i}\zeta_{i,j}-\sum_{j>i}\zeta_{i,j} \\
  x'_i\coloneq\lambda(i,\mathbb S)x_i+\zeta_i+\nabla x\,|\mathbb S|^{-1},~Y_i\coloneq x'_iG \\
  e_i^R\coloneq\uniform{\bits{\lambda}},~C_i^R\coloneq H(C,\texttt{"nonce"},i,R_i,e_i^R) \\
  \pcfor j\in\mathbb S\setminus\{i\};\pcdo \\
  \boldsymbol\beta_{j,i}\coloneq\uniform{\bits{L}},~\beta_{j,i}\coloneq\langle g_{j,i},\boldsymbol\beta_{j,i}\rangle \\
  (\rho^{\boldsymbol\beta_{j,i}}_{j,i},E_{j,i})\coloneq
    \call{SoftSpokenRecv}{A_{*,i,j},\boldsymbol\beta_{j,i}\mid c_{j,i},L} \\
  \pcsend{C_i^R,E_{j,i}}{j} \\
  \pcendfor{j} \\
  \mathrm{St}_i\coloneq\text{本轮本地变量、上下文及收发消息} \\
  \pcreturn
\end{protocol}

\begin{protocol}[
  name=PreSignRoundTwo,
  params={\mathrm{St}_i},
  returns={\mathrm{St}_i}
]
  d_i\coloneq H(C,\texttt{"nonce-view"},(C_j^R\mid j\in\mathrm{sort}(\mathbb S))) \\
  \pcfor j\in\mathbb S\setminus\{i\};\pcdo \\
  (\rho^0_{i,j},\rho^1_{i,j})\coloneq
    \call{SoftSpokenSend}{\delta_{*,j,i},A^*_{*,j,i},E_{i,j}\mid c_{i,j},L} \\
  (w_{i,j},\alpha^0_{i,j},Q_{i,j})\coloneq \\
  \quad\call{RvoleSend}{\rho^0_{i,j},\rho^1_{i,j}\mid c_{i,j},T_{i,j},L,2,1} \\
  y^a_{i,j}\coloneq-\sum_{\ell=0}^{L-1}g_{i,j,\ell}\alpha^0_{i,j,\ell,a},\quad a\in\{\mathrm r,\mathrm x\} \\
  \delta_{i,j}\coloneq(r_i,x'_i)-w_{i,j} \\
  \Gamma^{\mathrm r}_{i,j}\coloneq y^{\mathrm r}_{i,j}G,~\Gamma^{\mathrm x}_{i,j}\coloneq y^{\mathrm x}_{i,j}G \\
  \psi_{i,j}\coloneq\phi_i-\beta_{j,i} \\
  \pcsend{Q_{i,j},\delta_{i,j},R_i,e_i^R,Y_i,\Gamma^{\mathrm r}_{i,j},\Gamma^{\mathrm x}_{i,j},\psi_{i,j},d_i}{j} \\
  \pcendfor{j} \\
  \pcreturn
\end{protocol}

St 表示跨轮保留的本地状态，下一阶段继续使用并更新其中的变量；
上标 $\mathrm r,\mathrm x$ 分别对应 RVOLE 的功能负载 $0,1$。
SoftSpokenSend 失败时按 CheckP 处理其对端；RvoleSend 中的随机负载仅供本次使用。
接收第二轮后执行下列收尾，此处的 $\delta_{j,i}$ 是输入偏移，与树的打孔位置按索引形状区分。

\begin{protocol}[
  name=PreSignFinish,
  params={\mathrm{St}_i},
  returns={\mathrm{PreSignature}_i}
]
  \pcfor j\in\mathbb S\setminus\{i\};\pcdo \\
  \call{CheckP}{H(C,\texttt{"nonce"},j,R_j,e_j^R)=C_j^R,j} \\
  \call{CheckP}{d_j=d_i,\mathbb S\setminus\{i\}} \\
  t_{j,i}\coloneq\call{RvoleRecv}{\boldsymbol\beta_{j,i},\rho^{\boldsymbol\beta_{j,i}}_{j,i},Q_{j,i}\mid c_{j,i},T_{j,i},L,2,1} \\
  \pcfor a\in\{\mathrm r,\mathrm x\};\pcdo \\
  z^a_{j,i}\coloneq\sum_{\ell=0}^{L-1}g_{j,i,\ell}t_{j,i,\ell,a}+\beta_{j,i}\delta^a_{j,i} \\
  \pcendfor{a} \\
  \call{CheckP}{\beta_{j,i}R_j=\Gamma^{\mathrm r}_{j,i}+z^{\mathrm r}_{j,i}G,j} \\
  \call{CheckP}{\beta_{j,i}Y_j=\Gamma^{\mathrm x}_{j,i}+z^{\mathrm x}_{j,i}G,j} \\
  \pcendfor{j} \\
  \call{CheckP}{\sum_{j\in\mathbb S}Y_j=Y',\mathbb S\setminus\{i\}} \\
  R\coloneq\sum_{j\in\mathbb S}R_j,~\call{Check}{R\ne\mathcal O} \\
  r\coloneq R.x\bmod n,~\call{Check}{r\ne0} \\
  \Phi_i\coloneq\phi_i+\sum_{j\ne i}\psi_{j,i} \\
  V_i^a\coloneq\sum_{j\ne i}(y^a_{i,j}+z^a_{j,i}),\quad a\in\{\mathrm r,\mathrm x\} \\
  s_{1,i}\coloneq r_i\Phi_i+V_i^{\mathrm r} \\
  s^\circ_{0,i}\coloneq r(x'_i\Phi_i+V_i^{\mathrm x}) \\
  \mathrm{PreSignature}_i\coloneq(C,R,r,s_{1,i},s^\circ_{0,i},\phi_i) \\
  \pcreturn
\end{protocol}

RvoleRecv 失败时按 CheckP 处理发送方。
预签名连同上下文字段和一次性状态标记保存，其余临时矩阵与随机量销毁。
此时 $R$ 只是本方视图，完整点的一致性由第三轮检查。

\subsubsection{多方 Sign：第三轮}

采用 \href{https://eprint.iacr.org/2026/929.pdf}{2026/929，\S3}
的完整 nonce 点与签名符号检查。
读取预签名时先核对保存的消息、公钥、参与集合、参数与调用上下文；
发送签名份额之前原子地标记为已消费，成功、失败和崩溃恢复均不得再次使用。

\begin{protocol}[
  name=Sign,
  params={\mathrm{PreSignature}_i,m,Y',\mathbb S\mid C},
  returns={r,s},
  layout=multiline
]
  (C_i,R,r,s_{1,i},s^\circ_{0,i},\phi_i)\coloneq\mathrm{PreSignature}_i \\
  \call{Check}{C_i=C\ \land\ \mathrm{unused}(\mathrm{PreSignature}_i)} \\
  \call{Check}{R\ne\mathcal O\ \land\ r=R.x\bmod n\ \land\ r\ne0} \\
  s_{0,i}\coloneq s^\circ_{0,i}+m\phi_i \\
  \mathrm{Consume}(\mathrm{PreSignature}_i) \\
  \pcsend{C,R,s_{0,i},s_{1,i}}{\mathbb S\setminus\{i\}} \\
  \call{Check}{\forall j\ne i,~C_j=C\ \land\ R^{(j)}=R} \\
  D\coloneq\sum_{j\in\mathbb S}s_{1,j},~\call{Check}{D\ne0} \\
  s\coloneq D^{-1}\sum_{j\in\mathbb S}s_{0,j},~\call{Check}{s\ne0} \\
  \call{Check}{s^{-1}mG+s^{-1}rY'=R} \\
  \pcif\mathrm{signMode}=\texttt{low-S};\pcthen \\
  s\coloneq\min(s,n-s) \\
  \pcfi \\
  \call{Check}{\mathrm{EcdsaVerify}(m,Y',r,s)} \\
  \mathrm{Erase}(\mathrm{PreSignature}_i) \\
  \pcreturn
\end{protocol}

完整点比较包含纵坐标符号，符号等式在 low-S 规范化之前检查。
EcdsaVerify 验证 $1\le r,s<n$、公钥合法且非零，
并检查 $U=s^{-1}(mG+rY')\ne\mathcal O$ 及 $U.x\bmod n=r$。
所有失败出口也执行 Erase；Consume 的持久化先于网络发送。

\paragraph{现场零分享} 如选用 \href{https://eprint.iacr.org/2026/929.pdf}{2026/929，\S3}
的放宽零分享接口，令 $\mathrm{zeroMode}=\texttt{fresh}$，
Keygen 省去零分享种子的承诺、揭示和保存，其余初始化步骤照常执行。
PreSign 第一轮对每个 $i<j$，由 $i$ 采样 $a_{i,j}\gets\uniform{\mathbb Z_n}$ 并保密发送给 $j$。
第二轮开始时计算
\[
  \zeta_i=\sum_{j<i}a_{j,i}-\sum_{j>i}a_{i,j},\qquad
  x'_i=\lambda(i,\mathbb S)x_i+\zeta_i+\nabla x\,|\mathbb S|^{-1}.
\]
该模式将 $x'_i,Y_i$ 的计算移到第二轮生成 RVOLE 消息之前，轮数仍为三轮。
各方必须在 $C$ 中约定同一种模式，标量消息也绑定 $C$、双方编号与用途。

\subsubsection{两方独立预处理与两轮在线签名}

采用 \href{https://eprint.iacr.org/2026/976.pdf}{2026/976，Protocol 6.2}。
本节双方编号改为 $i\in\{0,1\}$，$j=1-i$，在线时提供有效且一致的加法私钥分片
$x'_0,x'_1$ 及公开分片 $Y_i=x'_iG$，共同公钥 $Y'=Y_0+Y_1$。
论文还要求兼容 Keygen 对恶意份额的提取；前述 Keygen 与该证明的对应仍需核实，
本节的代数编排不据此宣称完成了这项组合证明。

每份预处理使用唯一 $\mathtt{ppid}$，两个方向均独立初始化 Base OT 与 PPRF。
对每个方向，Receiver 调用 EndemicRecvInit，Sender 调用 EndemicSendInit，
双方同轮交换 $R^0,R^1$ 与 $V^0,V^1$，再分别调用 EndemicRecvFinish 和 EndemicSendFinish。
每 $k$ 个 OT 交给 PprfBuild，Sender 在下一轮发送各树的 $P$，Receiver 调用 PprfEval 校验。
会话及树编号、叶子保存和 $T^{\mathrm{init}}$ 的计算沿用 Keygen 的相应公式，
以 $\mathtt{ppid}$ 替换 $\mathtt{kgid}$；本次初始化不生成密钥份额。
这两轮初始化计入离线阶段，随后按下列协议执行两轮扩展与 RVOLE。

本节取 Variant II，$L=896,N_1=2,N_2=1$，有向实例号 $c_{i,j}$ 绑定 $\mathtt{ppid}$ 及双方角色。
$T_{i,j}$ 仍绑定该方向初始化和本批 $E_{i,j}$。
Nonce 揭示随机数采用本文的 $\lambda=256$ 比特口径。

\begin{protocol}[
  name=PreprocessTwoParty,
  params={A_{*,i,j},\delta_{*,j,i},A^*_{*,j,i},T^{\mathrm{init}}\mid\mathtt{ppid}},
  returns={\mathrm{State}_i},
  layout=multiline
]
  r_i\coloneq\uniform{\mathbb Z_n},~R_i\coloneq r_iG,~e_i^R\coloneq\uniform{\bits{\lambda}} \\
  C_i^R\coloneq H(\mathtt{ppid},\texttt{"nonce"},i,R_i,e_i^R) \\
  \boldsymbol\beta_{j,i}\coloneq\uniform{\bits{L}} \\
  (\rho^{\boldsymbol\beta_{j,i}}_{j,i},E_{j,i})\coloneq
    \call{SoftSpokenRecv}{A_{*,i,j},\boldsymbol\beta_{j,i}\mid c_{j,i},L} \\
  \pcsend{C_i^R,E_{j,i}}{j}\pccomment{离线扩展第一轮} \\
  (\rho^0_{i,j},\rho^1_{i,j})\coloneq
    \call{SoftSpokenSend}{\delta_{*,j,i},A^*_{*,j,i},E_{i,j}\mid c_{i,j},L} \\
  (w_{i,j},\alpha^0_{i,j},Q_{i,j})\coloneq \\
  \quad\call{RvoleSend}{\rho^0_{i,j},\rho^1_{i,j}\mid c_{i,j},T_{i,j},L,2,1} \\
  (\tilde r_i,\tilde x_i)\coloneq w_{i,j} \\
  g_{i,j}\coloneq\call{Gadget}{\mathrm{II},Q_{i,j}\mid c_{i,j},T_{i,j},L} \\
  a_i^b\coloneq-\sum_{\ell=0}^{L-1}g_{i,j,\ell}\alpha^0_{i,j,\ell,b},\quad b\in\{\mathrm r,\mathrm x\} \\
  \pcsend{Q_{i,j}}{j}\pccomment{离线扩展第二轮} \\
  t_{j,i}\coloneq\call{RvoleRecv}{\boldsymbol\beta_{j,i},\rho^{\boldsymbol\beta_{j,i}}_{j,i},Q_{j,i}\mid c_{j,i},T_{j,i},L,2,1} \\
  g_{j,i}\coloneq\call{Gadget}{\mathrm{II},Q_{j,i}\mid c_{j,i},T_{j,i},L} \\
  \beta_i\coloneq\langle g_{j,i},\boldsymbol\beta_{j,i}\rangle \\
  b_i^a\coloneq\sum_{\ell=0}^{L-1}g_{j,i,\ell}t_{j,i,\ell,a},\quad a\in\{\mathrm r,\mathrm x\} \\
  \mathrm{State}_i\coloneq(\beta_i,r_i,\tilde r_i,\tilde x_i,a_i^{\mathrm r},a_i^{\mathrm x},b_i^{\mathrm r},b_i^{\mathrm x}) \\
  \pcreturn
\end{protocol}

这里 $a,b$ 的下标表示份额持有人，因此
$a_i^{\mathrm r}+b_j^{\mathrm r}=\tilde r_i\beta_j$，
$a_i^{\mathrm x}+b_j^{\mathrm x}=\tilde x_i\beta_j$。
State 还保存对端承诺 $C_j^R$、本方揭示随机数 $e_i^R$、预处理编号、身份参数及待用标记。
成功压缩状态后销毁本批初始化秘密和临时矩阵；预处理失败也销毁整批材料。
此阶段不揭示 $R_i$，也不绑定实际签名密钥或消息。

在线时共同确定 $m,Y',Y_0,Y_1$，并检查 $Y'=Y_0+Y_1$ 及本方 $Y_i=x'_iG$。
在线上下文 $C_{\mathrm{on}}$ 绑定预处理编号、唯一签名号、双方身份、密钥标识、上述公开输入、参数及输出规则。
发送第一条在线消息之前，原子地将待用状态绑定 $C_{\mathrm{on}}$ 并标记已消费；
以后不能更换消息或密钥，取消、失败及崩溃恢复也不能将状态放回待用池。

\begin{protocol}[
  name=SignTwoParty,
  params={\mathrm{State}_i,x'_i,m,Y',Y_0,Y_1\mid C_{\mathrm{on}}},
  returns={r,s},
  layout=multiline
]
  \call{Check}{\mathrm{unused}(\mathrm{State}_i)\ \land\ Y'=Y_0+Y_1\ \land\ Y_i=x'_iG} \\
  \mathrm{Consume}(\mathrm{State}_i,C_{\mathrm{on}}) \\
  \gamma_i^{\mathrm r}\coloneq r_i-\tilde r_i,~\gamma_i^{\mathrm x}\coloneq x'_i-\tilde x_i \\
  \Gamma_i^{\mathrm r}\coloneq a_i^{\mathrm r}G,~\Gamma_i^{\mathrm x}\coloneq a_i^{\mathrm x}G \\
  R_i\coloneq r_iG \\
  \pcsend{C_{\mathrm{on}},\gamma_i^{\mathrm r},\gamma_i^{\mathrm x},\Gamma_i^{\mathrm r},\Gamma_i^{\mathrm x},R_i,e_i^R}{j}
    \pccomment{在线第一轮} \\
  \call{Check}{C_{\mathrm{on}}^{(j)}=C_{\mathrm{on}}} \\
  \call{Check}{H(\mathtt{ppid},\texttt{"nonce"},j,R_j,e_j^R)=C_j^R} \\
  z_i^{\mathrm r}\coloneq b_i^{\mathrm r}+\beta_i\gamma_j^{\mathrm r},~
    z_i^{\mathrm x}\coloneq b_i^{\mathrm x}+\beta_i\gamma_j^{\mathrm x} \\
  \call{Check}{z_i^{\mathrm r}G=\beta_iR_j-\Gamma_j^{\mathrm r}} \\
  \call{Check}{z_i^{\mathrm x}G=\beta_iY_j-\Gamma_j^{\mathrm x}} \\
  R\coloneq R_0+R_1,~\call{Check}{R\ne\mathcal O} \\
  r\coloneq R.x\bmod n,~\call{Check}{r\ne0} \\
  u_i\coloneq r_i\beta_i+a_i^{\mathrm r}+z_i^{\mathrm r} \\
  v_i\coloneq x'_i\beta_i+a_i^{\mathrm x}+z_i^{\mathrm x} \\
  w_i\coloneq m\beta_i+rv_i \\
  \pcsend{C_{\mathrm{on}},u_i,w_i}{j}\pccomment{在线第二轮} \\
  \call{Check}{C_{\mathrm{on}}^{(j)}=C_{\mathrm{on}}} \\
  D\coloneq u_0+u_1,~\call{Check}{D\ne0} \\
  s\coloneq(w_0+w_1)D^{-1},~\call{Check}{s\ne0} \\
  \call{Check}{\mathrm{EcdsaVerify}(m,Y',r,s)} \\
  \pcif\mathrm{signMode}=\texttt{low-S};\pcthen \\
  s\coloneq\min(s,n-s) \\
  \pcfi \\
  \mathrm{Erase}(\mathrm{State}_i) \\
  \pcreturn
\end{protocol}

所有失败出口均销毁该批一次性状态。
两方流程直接以 $\beta_i$ 作为盲化分片，满足
\[
  \begin{aligned}
    u_0+u_1&=(r_0+r_1)(\beta_0+\beta_1),\\
    v_0+v_1&=(x'_0+x'_1)(\beta_0+\beta_1).
  \end{aligned}
\]
因此其两轮在线阶段使用自己的消息格式，按 Protocol 6.2 验签后输出。

\end{document}
